%======================================================================
\section{The mixing problem as a marked split--merge chain: symmetry,
factorization, and the exact rate from below}
\label{sec:mixing}

This section records the session-8 progress on
Problem~\ref{prob:mixing}.  The fresh-chord averaging operator $P$ of
\S\ref{sec:factor} is identified \emph{exactly} with one step of the
uniform split--merge (coagulation--fragmentation) dynamics on measured
partitions of the circle, carrying two marked points
(Lemma~\ref{lem:splitmerge}); the chain admits an exact sign-reversing
involution (Theorem~\ref{thm:iota}) which factors the fresh-averaged
sign as $M_d=\varepsilon_0\cdot\Phi_d(\text{geometry})$
(Corollary~\ref{cor:quotient}) and removes the label from the problem;
the stationary law of the relevant ``gap'' coordinate is computed
exactly (linear density; Proposition~\ref{prop:gaplawstat}); and these
combine into a \emph{proved} lower bound
$\bar\psi(d)\ge c/d$ (Proposition~\ref{prop:lower}): the numerically
measured rate $\psi\approx1.0/d$ is sharp, and $c=1/2$ is the true
exponent of Problem~\ref{prob:mixing} if the matching upper bound
holds.  The upper bound is reduced to two sharply scoped sub-problems
(Problems~\ref{prob:gaplaw} and~\ref{prob:survival}), each with the
mechanism confirmed numerically, plus a spectral route via
reversibility (\S\ref{sec:spectral}).  \emph{Session 9
(\S\ref{sec:s9}) shows the first of the two is unnecessary over the
horizon $j\le8d$ relevant to Theorem~\ref{thm:factor}: any
$\alpha>2/9$ in Problem~\ref{prob:survival} closes the iid model.
Session 10 (\S\ref{sec:s10}) shows the environment must enter that
problem through recharge, proves that the block count on the horizon is
$O(\log d)$ (the Harer--Zagier law of the number of curves), so the
environment is healthy where it is needed, and reduces
Problem~\ref{prob:survival} to two conditions (R1)--(R2) on the
log-macroscopic post-flip chain via an exact Markov-renewal identity.}

Throughout, the ensemble is that of \S\ref{sec:factor}: $u,v$ fixed on
the circle of unit circumference ($u$ at $0$, $v$ at $1/2$ in the
code), chords $=$ iid uniform endpoint pairs.

%----------------------------------------------------------------------
\subsection{Identification with the uniform split--merge chain}
\label{sec:smident}

For a finite chord set $W$, the cycles of $\gamma\tau_W$ trace out a
partition of the circle (minus the endpoint set) into finite unions of
arcs (``curves''); each curve carries its Lebesgue measure, and its
internal cyclic order (the order in which $\gamma\tau_W$ traverses its
arcs) transports Lebesgue measure to a uniform measure along the
curve.  The \emph{marked state} is
\[
X \;=\;\Bigl(\{\text{curve masses}\}, \ \text{curves of }u,v,\
\begin{cases}
(x,y) & u\sim v:\ \text{masses of the two }u\to v\text{ curve arcs},\\
(p,q) & u\not\sim v:\ \text{masses of the curves through }u,\ v,
\end{cases}\Bigr)
\]
and $\varepsilon(W)=+1$ iff $u\sim v$.

\begin{lemma}[fresh chord $=$ one split--merge move]
\label{lem:splitmerge}
Adding one fresh chord (iid uniform endpoints $c_1,c_2$) maps the
marked state by the following rules, which depend on the current state
only.  Write $B(c)$ for the curve containing $c$; positions of $c_i$
along their curves are independent uniform.
\begin{enumerate}[leftmargin=2em]
\item $B(c_1)=B(c_2)=B$ (probability $|B|^2$): $B$ splits into the two
arcs of $B$ between the cuts.  Sub-cases:
 (a) $B$ unmarked, or marked by exactly one of $u,v$ (mass $p$): the
 mark keeps the piece of mass $p-|t_1-t_2|$, $t_i\sim U(0,p)$ iid; the
 complementary piece of mass $|t_1-t_2|$ joins the environment.
 (b) $B$ the common block, sides $(x,y)$: with the cycle parametrised
 $u\to v$ ($x$-side) then $v\to u$ ($y$-side), cuts uniform on
 $(0,x{+}y)$: both cuts on one side shrinks that side as in (a);
 one cut per side (probability $2xy/(x{+}y)^2$ given the split)
 \emph{separates} $u,v$: the new blocks have masses
 $q'=t_2-t_1$ ($v$'s) and $p'=x+y-q'$ ($u$'s), i.e.\
 $(p',q')\overset{d}{=}(\alpha+y-\beta,\ x-\alpha+\beta)$,
 $\alpha\sim U(0,x)$, $\beta\sim U(0,y)$.
\item $B(c_1)\ne B(c_2)$: the two curves merge (cut-and-join at
$c_1,c_2$).  Sub-cases: (a) at most one is marked, or one holds $u$
and the other neither: masses add, marks carried along.  (b) $u$'s and
$v$'s blocks (masses $p,q$) merge: $u,v$ become joined with sides
$(x,y)=(\alpha+\beta,\ p+q-\alpha-\beta)$, $\alpha\sim U(0,p)$,
$\beta\sim U(0,q)$.  (c) the common block (sides $x,y$) merges with an
environment block of mass $z$: the insertion point is on the $x$-side
with probability $x/(x{+}y)$, and that side gains $z$.
\end{enumerate}
In particular $P$ (add one fresh chord and average) is the transition
operator of this \emph{marked uniform split--merge chain}, the
continuum limit of the cycle structure of the random-transposition
walk (cf.\ \cite{Sch05}), and
$M_d=P^d\varepsilon$ is its $d$-step averaged together-indicator sign.
\end{lemma}

\begin{proof}
Adding the chord $z=(c_1,c_2)$ replaces $\gamma\tau_W$ (computed on
the endpoint set enlarged by $c_1,c_2$; refining $\gamma$ at
non-endpoint positions does not change the curve partition) by
$\gamma\tau_W\,(c_1c_2)$.  Right multiplication by a transposition
whose entries lie on one cycle splits it into the two arcs between
them, and merges two cycles otherwise --- with the traversal orders
spliced at the cuts; this is the standard surgery computation.  The
probabilistic statements: $c_1,c_2$ are iid Lebesgue, the curves
partition the circle, and Lebesgue restricted to a curve is uniform
along it; all sub-case laws then read off from the traversal order,
e.g.\ in 1(b) the cut positions along the cycle are iid $U(0,x{+}y)$
and $u,v$ separate iff the cuts straddle both sides.
\end{proof}

\emph{Status: PROVED.  VERIFIED end-to-end
(\texttt{s8\_splitmerge.c}): the continuum chain reproduces the
discrete-surgery values of $\psi(j,d)$ (fs$=1$ unbiased runs;
$\psi(16,4)$: $0.1341(22)$ vs $0.1305(22)$;
$\psi(64,16)$: $0.0391(26)$ vs $0.0395(26)$), the one-point function
($f(4)$: $+0.0232(4)$ vs $+0.0237(5)$), and the
Poisson--Dirichlet sanity values at stationarity
($\E\sum p_i^2=0.4979$ vs $\mathrm{PD}(1)$'s $1/2$;
$\E[\text{common-block mass}\mid u\sim v]=0.6685$ vs $2/3$, see
Proposition~\ref{prop:gaplawstat}).  NOTE: repeat runs show the
session-7 $\psi$ table of \S\ref{sec:factor} carries sampling scatter
$\approx\pm0.005$ (its inner-average standard errors were optimistic);
its fitted exponent is unaffected.}

%----------------------------------------------------------------------
\subsection{The sign-reversing involution and the quotient chain}
\label{sec:iota}

\begin{theorem}[label-flip symmetry]\label{thm:iota}
Let $\iota$ be the involution of marked states that swaps the label
$\{u\sim v\}\leftrightarrow\{u\not\sim v\}$ while keeping the
environment and the ordered geometry pair fixed:
\[
\iota:\ \bigl(\text{together},(x,y),\mathrm{env}\bigr)
\;\longleftrightarrow\;
\bigl(\text{apart},(p,q)=(x,y),\mathrm{env}\bigr).
\]
Then $\iota$ commutes with the transition kernel of the marked
split--merge chain, and $\varepsilon\circ\iota=-\varepsilon$.
Consequently $M_d\circ\iota=-M_d$ for every $d$.
\end{theorem}

\begin{proof}
Match the transition laws of Lemma~\ref{lem:splitmerge} case by case
under the correspondence $x\leftrightarrow p$ (the $u$-coordinate),
$y\leftrightarrow q$ ($v$-coordinate).  The total marked mass is
$x+y=p+q$ in both states, so the probabilities of the events ``both
cuts in marked material'', ``one cut in marked material, one in a
given environment block'', ``both in environment'' agree, as does the
finer classification by $u$- vs $v$-material ($x^2/(x{+}y)^2$ etc.,
against $p^2/(p{+}q)^2$ etc.).  Given the event: both cuts in
$u$-material shrinks $x$ (resp.\ $p$) by $|t_1-t_2|$ with the same
law; a cut pair straddling $u$- and $v$-material is the
\emph{separation} in the together state and the \emph{join} in the
apart state --- i.e.\ it flips the label in both --- and the outgoing
geometry laws agree:
$(\alpha+y-\beta,\,x-\alpha+\beta)
\overset{d}{=}(\alpha+\beta',\,x+y-\alpha-\beta')$ jointly, by
$\beta\mapsto y-\beta'$; a single cut in $u$-material merges an
environment block $z$ into the $u$-coordinate in both states; and
environment-only moves are identical.  This exhausts all cases, so
$K(\iota s,\iota\,\cdot)=\iota_*K(s,\cdot)$; $\varepsilon$ flips by
definition of $\iota$, and $M_d=P^d\varepsilon$ inherits oddness.
\end{proof}

\emph{Status: PROVED (case check as above).  VERIFIED
(\texttt{s8\_splitmerge iota}): from mirrored starts
$(\mathrm{tog},(x,y),z)$ and $(\mathrm{apart},(x,y),z)$, measured
$P_A(u\sim v)+P_B(u\sim v)=1$ within Monte Carlo error at all tested
$t\le32$ and geometries, with identical $u$-coordinate marginals.}

\begin{corollary}[quotient factorization; the label carries all the
sign]\label{cor:quotient}
Let $\Gamma_t=(a_t,b_t;\mathrm{env}_t)$ be the quotient of the marked
chain by $\iota$ (the geometry with the label forgotten;
$a\leftrightarrow x$ or $p$, $b\leftrightarrow y$ or $q$).  Then:
\begin{enumerate}[leftmargin=2em]
\item $\Gamma_t$ is Markov, with the transition laws of
Lemma~\ref{lem:splitmerge} read on $(a,b;\mathrm{env})$; the
\emph{flip moves} (cut pair straddling $a$- and $b$-material) occur
with probability $2a_tb_t$ per step and update
$(a,b)\to(\alpha+\beta,\,a+b-\alpha-\beta)$.
\item The label evolves deterministically along the quotient path:
$\varepsilon_d=\varepsilon_0\cdot(-1)^{N_d}$ with $N_d$ the number of
flip moves, a functional of the quotient path alone.
\item Hence, with
$\Phi_d(\Gamma):=\E_\Gamma\bigl[(-1)^{N_d}\bigr]$,
\[
M_d(C)\;=\;\varepsilon(C)\cdot\Phi_d\bigl(\Gamma(C)\bigr),
\qquad
\psi(j,d)\;=\;\E_{\mu_j}\bigl[\Phi_d(\Gamma_0)^2\bigr].
\]
\item Every stationary law of the marked chain is $\iota$-invariant
wherever uniqueness holds on its support; in particular the label is
a fair coin given the geometry at stationarity.
\end{enumerate}
\end{corollary}

\begin{proof}
(1),(2): the flip moves are exactly the label-changing moves in both
sectors (Theorem~\ref{thm:iota}), and their effect on the geometry has
the stated common law; all other moves preserve the label.  (Flip
moves are a.s.\ identifiable on the quotient path: they are the only
moves that change both marked coordinates.)  (3):
condition on the quotient path.  (4): $\iota$ pushes any stationary
law to a stationary law with the same quotient marginal.
\end{proof}

\emph{Status: PROVED.  Problem~\ref{prob:mixing} is thereby freed of
the sign structure: it is now a statement about the \emph{parity}
functional $\Phi_d=\E[(-1)^{N_d}]$ of an autonomous
positive-recurrent chain.  Consistency check: at stationarity the
label-fairness (4) predicts $P(u\sim v)=1/2$; measured $0.5065(29)$ at
burn-in $8192$.}

\begin{remark}[renewal structure at flip times]\label{rem:renewal}
Let $\tau$ be the first flip move.  Strong Markov gives the exact
renewal identity
\[
\Phi_d(\Gamma)\;=\;\Pr_\Gamma[\tau>d]\;-\;
\E_\Gamma\bigl[\Phi_{d-\tau}(\Gamma_\tau);\ \tau\le d\bigr].
\]
The flip clock runs at rate $2a_tb_t$; after a flip the $u$-coordinate
is \emph{recharged} ($a'=\alpha+\beta$ is macroscopic unless both
$a,b$ were small).  The correlation-survival mechanism is therefore
gap-trapping: $\Pr[\tau>d]$ is dominated by trajectories on which
$g_t=\min(a_t,b_t)$ stays $O(1/d)$, since an untouched small side
freezes ($\Pr[\text{the }g\text{-material is hit in }d\text{ steps}]
\le 4gd$), while a touched one is typically recharged to macroscopic
size by a size-biased merge.
\end{remark}

%----------------------------------------------------------------------
\subsection{The stationary gap law, and the exact rate from below}
\label{sec:gaplaw}

\begin{proposition}[stationary gap law]\label{prop:gaplawstat}
Let $\pi$ be the stationary law of the marked chain (the $n\to\infty$
limit of the cycle structure of a uniform permutation of $n$ circle
points with two marked points; top of the size spectrum
$\mathrm{PD}(1)$).  Under $\pi$: $\Pr[u\sim v]=\tfrac12$; given
$u\sim v$ the common-block mass $s$ has density $2s$ on $(0,1)$ (so
$\E[s\mid u\sim v]=2/3$) and $x\mid s\sim U(0,s)$; given $u\not\sim v$
the pair $(p,q)$ is uniform on the triangle $\{p,q>0,\ p+q<1\}$
(so the joint density of $(\text{apart},p,q)$ is $1$ there, mass
$\tfrac12$).  Consequently, with
$g=\min(x,y)$ resp.\ $\min(p,q)$,
\[
\Pr_\pi[g\le\delta]\;=\;4\delta+O(\delta^2)\qquad(\delta\to0),
\]
and the density of $g$ at $0^+$ is $4$.
\end{proposition}

\begin{proof}
At $n$ points: for uniform $\sigma$, $\Pr[u,v$ in a common $m$-cycle$]
=(m-1)/\bigl(n(n-1)\bigr)$ (choose and order the cycle), summing to
$1/2$; given the cycle, the $u\to v$ distance is uniform on
$\{1,\dots,m-1\}$; if apart,
$\Pr[|C_u|=m_1,v\notin C_u]=\tfrac1n\cdot\tfrac{n-m_1}{n-1}$ and
given this $|C_v|$ is uniform on $\{1,\dots,n-m_1\}$, so the joint
law of $(p,q)=(m_1,m_2)/n$ converges to
$(1-p)\,dp\cdot\tfrac{dq}{1-p}=dp\,dq$ on the triangle.  Pass to the
limit ($m/n\to s$).  The
gap: given together,
$\Pr[\min(x,y)\le\delta]=\int_0^1 2s\min(1,2\delta/s)\,ds
=4\delta-4\delta^2$; given apart,
$\Pr[\min(p,q)\le\delta]=4\delta+O(\delta^2)$ likewise; average the
two sectors.
\end{proof}

\emph{Status: PROVED at finite $n$ with a routine limit exchange
(the marked chain is the lumping of the transposition walk, see
\S\ref{sec:spectral}; over any fixed horizon the kernels converge).
VERIFIED: measured dyadic masses at stationarity
($\Pr[g\in(2^{-6},2^{-5})]$: together $0.0297$ vs predicted $0.0312$,
apart $0.0288$ vs $0.0312$; lower bins consistent with density $4$
throughout), and the label is a fair coin in every gap bin.}

\begin{proposition}[the mixing rate is at most $1/d$: lower bound]
\label{prop:lower}
There are $c_0>0$, $d_0$ such that the stationary-start mixing
functional obeys
\[
\psi_\pi(d)\;:=\;\E_\pi\bigl[M_d^2\bigr]\;\ge\;\frac{c_0}{d},
\qquad d\ge d_0 .
\]
Hence $\bar\psi(d)\ge c_0/d$ along any initial family
$\{\mu_j\}$ whose gap law dominates a multiple of the stationary one
(numerically true for the iid-chord $\mu_j$ already at $j=16$), and
no exponent $c>1/2$ is attainable in Problem~\ref{prob:mixing}.
\end{proposition}

\begin{proof}
Fix $K$ and let
$G=\{g_0\le1/(Kd)\}$.  On $G$, the $g$-material (the short side if
together, the small block if apart) is hit by no cut during $d$ steps
except with probability $\le 4g_0d\le4/K$: each step's two cuts avoid
it with probability $\ge1-4g_0$ while it remains unhit, and while
unhit its mass is frozen.  If it is never hit there is no flip move
(a flip requires a cut in both $a$- and $b$-material), so
$\varepsilon_d=\varepsilon_0$ and, for the conditional expectation,
$|M_d|=|\E[\varepsilon_d\mid X_0]|\ge1-8/K$ on $G$.  Therefore
$\psi_\pi(d)\ge\E[M_d^2;G]\ge(1-8/K)^2\,\pi(G)$, and
$\pi(G)\ge c_1/(Kd)$ by Proposition~\ref{prop:gaplawstat}.  Take
$K=16$.
\end{proof}

\emph{Status: PROVED (given Proposition~\ref{prop:gaplawstat}).  This
settles the sharpness question left open in \S\ref{sec:factor}: the
measured $\psi\approx0.8/d$ is the true order, the budget margin of
Corollary~\ref{cor:psienough} is exactly $c=1/2$ vs $c>1/11$
(factor $5.5$ in the exponent), and no route can do better than
$1/d$ --- in particular any proof of Problem~\ref{prob:mixing} must
tolerate the gap-trapped trajectories, they are really there.}

%----------------------------------------------------------------------
\subsection{What remains for the upper bound}
\label{sec:upper}

The decomposition
$\psi(j,d)=\E[\Phi_d^2;\,g_0<\varepsilon]+\E[\Phi_d^2;\,g_0\ge
\varepsilon]$ with $\varepsilon\asymp1/d$, together with
$|\Phi_d|\le1$, reduces the upper bound to two independent
statements.

\begin{problem}[uniform gap law]\label{prob:gaplaw}
Show $\Pr_{\mu_j}[g_0\le\varepsilon]\le C\varepsilon^{\kappa}$ (some
$\kappa\in(0,1]$, uniformly in $j$) for the iid-chord initial laws
$\mu_j$ --- equivalently for the chain run $j$ steps from the
one-block start.
\end{problem}

\emph{Partial progress: (i) at stationarity this is
Proposition~\ref{prop:gaplawstat} with $\kappa=1$; (ii) the per-step
\emph{creation} flux into $\{g\le\varepsilon\}$ has the universal
bound $\Pr[\text{step creates }g'\in d\varepsilon]\le
C\varepsilon\,d\varepsilon$ for every current state (each channel ---
shrink of a marked coordinate, flip recharge landing small --- has
outgoing density $\le2\varepsilon$ at $\varepsilon$: e.g.\ a split of
the $u$-block of mass $p$ leaves $u$ with $p-|t_1-t_2|$, of density
$2\varepsilon'/p^2\cdot p^2$ at $\varepsilon'$); (iii) escape from
$\{g\le\varepsilon\}$ happens at rate $\ge2g(1-g)\cdot
\Pr[\text{merge partner macroscopic}]$, so a renewal bound gives
$\kappa=1$ \emph{given} a uniform-in-time ``no-dust'' estimate (the
environment retains a macroscopic block); it is the no-dust step that
is open, and only it.  Numerically the dyadic gap masses under
$\mu_j$ match the stationary density already at $j=16$.}

\begin{lemma}[recharge estimates]\label{lem:recharge}
For the quotient chain: (i) flip moves preserve the marked total
$s=a+b$; (ii) a flip from $(a,b)$ lands at gap
$g'=\min(\alpha+\beta,\,s-\alpha-\beta)$ with
$\Pr[g'\le\eta]\le\eta^2/(ab)$ for $\eta\le\min(a,b)$; (iii) from
\emph{any} state, each move channel sends a marked coordinate into
$(0,\eta]$ from above with probability $\le\eta^2$ per step (e.g.\
the shrink channel: both cuts in the $a$-material has probability
$a^2$, and $|t_1-t_2|\in[a-\eta,a)$ has conditional probability
$\le\eta^2/a^2$); so the total creation flux into
$\{g\le\eta\}$ is $\le4\eta^2$ per step.
\end{lemma}

\begin{proof}
(i) is immediate from the flip law.  (ii): $\alpha+\beta\le\eta$ has
probability $\eta^2/(2ab)$ ($\alpha\sim U(0,a)$, $\beta\sim U(0,b)$
independent), likewise the other end.  (iii): direct integration of
the stated densities, channel by channel (shrink of $a$, shrink of
$b$, and the two flip ends).
\end{proof}

\emph{Status: PROVED.  Together with Remark~\ref{rem:renewal} this
pins the intended endgame: from a good start the first-flip time
$\tau$ has tail $\asymp1/t$ (trapping is the only way to stall, and
Lemma~\ref{lem:recharge}(iii) prices its creation), each flip
recharges the gap with quadratic safety (ii), so the flips form an
approximate renewal process with infinite-mean, regularly-varying
inter-arrival tail $\sim c/t$; for such renewals the alternating sum
$\E[(-1)^{N_d}]$ collapses to $\Theta(\text{tail}(d))=\Theta(1/d)$
(the signed renewal measure $\sum_k(-1)^kG^{*k}$ has total mass
$\tfrac12$).  Two genuine gaps remain: the inter-flip times are
state-dependent rather than iid, and the escape-rate lower bound
$\ge c\,g$ behind the $1/t$ tail needs the environment to offer a
macroscopic merge partner --- a ``no-dust'' estimate; at stationarity
dust configurations have probability $\approx$ the Dickman mass
$\rho(8)\sim10^{-8}$, small but not $o(1)$ in $d$, so the assembly
must let dust episodes delay rather than veto escape (dust coarsens
under merges).  Note also the second trap coordinate: the marked total
$s$ itself can be small (then flips are rare no matter the gap); but
$\Pr_\pi[s\le\delta]=O(\delta^2)$ in both sectors (the together mass
$s$ has density $2s$, and $(p,q)$ is uniform on the triangle), so the
$s$-trap contributes only $O(1/d^2)$ and rides along harmlessly in
the renewal bookkeeping.}

\begin{problem}[parity decay off small gaps]\label{prob:survival}
Show $\E_{\mu_j}\bigl[\Phi_d(\Gamma_0)^2;\ g_0\ge\varepsilon\bigr]\le
C(\varepsilon d)^{-2\alpha}$ for some $\alpha>0$ (uniformly in $j$),
i.e.\ trajectories started with both marked coordinates
$\ge\varepsilon$ lose the parity memory at a polynomial rate in
$\varepsilon d$.
\end{problem}

Any pair $(\kappa,\alpha)$ solves Problem~\ref{prob:mixing}: from
$\psi\le C\bigl(\varepsilon^\kappa+(\varepsilon d)^{-2\alpha}\bigr)$,
optimising at $\varepsilon=d^{-2\alpha/(\kappa+2\alpha)}$ gives
$\psi\le Cd^{-2c'}$ with $c'=\alpha\kappa/(\kappa+2\alpha)$; the
budget needs only $c'>1/11$, e.g.\ $(\kappa,\alpha)=(1,1/2)$ gives
$c'=1/4$, and any $\kappa=\alpha>3/10$ still closes
($c'=\kappa/3>1/11$ there).  \emph{The mechanism behind
Problem~\ref{prob:survival} is confirmed numerically: binning
$\psi$-contributions by $g_0$ (\texttt{s8\_splitmerge psibar}) shows
$\E[\Phi_d^2\mid g_0]$ is a scaling function of $g_0d$ alone ---
$\approx1$ for $g_0d\lesssim1/4$, $\lesssim0.05$ for $g_0d\gtrsim4$
--- so the conditional decay is fast; what is missing is a proof
device for the parity cancellation (the renewal identity of
Remark~\ref{rem:renewal} plus recharge is the intended route: after
the first flip the state is macroscopic w.h.p., flips then recur at
rate $\Theta(1)$, and the alternating sum cancels to the heavy-tail
term $\Pr[\tau>d]$).}

%----------------------------------------------------------------------
\subsection{The spectral route}\label{sec:spectral}

The discrete precursor of the marked chain --- right multiplication
of $\sigma_t\in S_n$ by uniform transpositions, marked cycle
structure observed --- is a strongly lumpable function of a chain
that is reversible with respect to the uniform measure; the lumped
chain is therefore reversible with respect to its stationary law, and
this persists in the $n\to\infty$ limit chainwise (the continuum
statement deserves its own write-up; nothing below is used elsewhere).
Granting it: $\psi_\pi(d)=\int\lambda^{2d}\,d\nu_\varepsilon(\lambda)$
for the spectral measure $\nu_\varepsilon$ of $\varepsilon$ under the
self-adjoint $P$, and Problem~\ref{prob:mixing} (stationary start)
becomes: \emph{$\nu_\varepsilon$ has at most polynomial mass near
$|\lambda|=1$}, with Proposition~\ref{prop:lower} asserting at least
linear mass near $\lambda=1$.  Two structural facts make this
attractive.  (i) By Theorem~\ref{thm:iota} the sign lives in the
$\iota$-odd sector, where the Dirichlet form has a \emph{massive}
term on flip moves: for $f=\varepsilon\cdot F(\Gamma)$,
\[
\mathcal{E}(f,f)=\tfrac12\,\E_\pi\Bigl[(F-F')^2\,;\,\text{non-flip}
\Bigr]+\tfrac12\,\E_\pi\Bigl[(F+F')^2\,;\,\text{flip}\Bigr],
\]
so a low-energy odd function must be small wherever the flip rate
$2ab$ is not, i.e.\ must localise on $\{g\lesssim\delta\}$, whose
mass is $\asymp\delta$ (Proposition~\ref{prop:gaplawstat}): this is
precisely the geometric picture of a linear spectral density, and
trial functions $F=\mathbf{1}\{g\le\delta\}$ realise it (Rayleigh
quotient $\asymp\delta$ at mass $\asymp\delta$), reproving
Proposition~\ref{prop:lower} spectrally.  (ii) The converse direction
--- \emph{every} odd $F$ with
$\mathcal{E}\le\delta\|F\|^2$ has
$|\langle F,\varepsilon\rangle|^2\le C\delta\|F\|^2$ --- is a
weighted Hardy-type inequality in the gap coordinate and is the
cleanest formulation of the missing upper bound.  A caveat for
whoever takes this up: the chain is $2$-periodic
($(-1)^{\#\mathrm{blocks}}$ flips deterministically, an exact
$\lambda=-1$ eigenfunction, and it is $\iota$-odd since $\iota$
changes the block count by one), so the top-of-spectrum analysis must
be run on $P^2$ or per parity sector; $\psi$ decays, so
$\langle\varepsilon,(-1)^{\#\mathrm{blocks}}\rangle_\pi=0$, which is
consistent with phase symmetry of $\pi$, but intermediate
inequalities must not silently assume aperiodicity.

%----------------------------------------------------------------------
\subsection{Transfer to the orbit measure: the mixture route}
\label{sec:transfer}

The factorization proof of Theorem~\ref{thm:factor} uses exactly one
probabilistic input --- chords attached to disjoint index sets are
conditionally independent given the shared ones --- and the session-7
caution stands: this fails per-instance.  The split--merge picture
suggests the correct weakening.  The orbit chord positions
$(Z_1,\dots,Z_k)$ are exchangeable given the frame (the cycles of
$\pi_P$, the points $u,v$); a finitely exchangeable vector drawn from
a population of $N$ position slots is a mixture of iid$(\theta)$
vectors up to total variation $O(k^2/N)$ (Diaconis--Freedman), with
$\theta$ the (random) empirical position law of the instance, and the
cell-disjointness of the intercalate family is a pairwise hard-core
correction of the same $O(k^2/N)$ order.  Two observations make the
mixture usable:
\begin{enumerate}[leftmargin=2em]
\item Conditionally on $\theta$, Theorem~\ref{thm:factor} applies
verbatim (its proof needs only iid given $(u,v,\theta)$), giving
$4\,\E[\mathrm{bias}^2]\le
\E_\theta\bigl[\max_{j,d\ge k/8}\psi^\theta(j,d)\bigr]
+2e^{-k/32}+O(k^2/N)$.
\item The entire analysis of
\S\S\ref{sec:smident}--\ref{sec:gaplaw} is \emph{measure-covariant}:
for chords with iid $\theta$-distributed endpoints, the induced chain
is the split--merge chain in which every rate is read off the
$\theta$-masses of curves and sides (cuts land in a curve with
probability its $\theta$-mass, uniformly in its
$\theta$-parametrisation).  In particular Lemma~\ref{lem:splitmerge},
Theorem~\ref{thm:iota}, Corollary~\ref{cor:quotient} and
Lemma~\ref{lem:recharge} hold verbatim with ``mass'' $=$
$\theta$-mass, provided $\theta$ is atomless; atoms of mass $m$ act
as permanent traps of that scale and cap the decay at
$\psi^\theta\gtrsim\max_i\theta(\{i\})$, which for orbit frames
($N\asymp n$ near-uniform slots) is $O(1/n)$ --- far below the
$k^{-2c}$ budget, $k=\Theta(\log^{11}n)$.
\end{enumerate}
So the transfer problem 7a(i) is reduced to (a) a quantitative
exchangeable-to-mixture step for the greedy maximal family under the
orbit measure (the hard-core correction; KSS-type concentration
should control the conflict density), and (b) the
$\theta$-uniform version of the mixing bound, i.e.\
Problem~\ref{prob:survival} with constants depending only on
$\max_i\theta(\{i\})$ and not otherwise on $\theta$
(Proposition~\ref{prop:shorthorizon} is already $\theta$-covariant:
Lemma~\ref{lem:recharge}(iii) reads verbatim in $\theta$-mass for
atomless $\theta$).  Empirically the conclusion already holds on real squares:
binning the session-7 per-instance data
(\texttt{s7\_inst\{30,50,100\}.csv}) by $k$ gives
$\E[\mathrm{bias}^2\mid k]\cdot k=0.6$--$0.75$, flat across
$k=8$--$26$ and $n=30$--$100$ (and \texttt{s8\_inst200.csv} extends
this, see \S\ref{sec:mixnum}): the $1/k$ law transfers with a
constant $\approx2.2\times$ the iid one, no drift.

\emph{Status: reduction formulated; both steps OPEN.  This replaces
the vaguer options ($\alpha$)/($\beta$) of the session-7 list.}

%----------------------------------------------------------------------
\subsection{Numerics at large $d$}\label{sec:mixnum}

All session-8 runs, continuum chain unless stated
(\texttt{s8\_psibar.log}, \texttt{s8\_4eb2\_big.log},
\texttt{s8\_onept.log}); COMPUTED.

\emph{(a) The stationary rate is $1/d$ on the nose.}  Stationary-start
(burn-in $16384$) values of $\psi_\pi(d)\cdot d$:
$0.96$ ($d{=}64$), $1.06$ ($d{=}128$), $1.11$ ($d{=}256$),
$0.93$ ($d{=}512$), $0.87$ ($d{=}1024$; errors grow from $\pm0.05$
to $\pm0.25$ along the list) --- constant $\approx1.0$, fitted
exponent $1.03\pm0.09$ over $d=64$--$1024$ (and $1.05$ on the
session-7 $j$-table route), no logarithmic growth, matching
Proposition~\ref{prop:lower} from below and pinning the conjectured
upper bound at $c=1/2$ exactly.

\emph{(b) The scaling collapse.}  Binning the $\psi$-contribution by
the initial gap: $\E[\Phi_d^2\mid g_0]$ is a function of $u=g_0d$
alone across $d=64,128,512$ --- the dyadic profile
$\varphi(u)\approx0.06,\,0.23,\,0.49,\,0.69,\,0.82,\,0.9,\,{\to}1$ at
$u\approx0.7,\,0.35,\,0.18,\,0.09,\,0.04,\,0.02,\dots$, and
$\varphi\le0.02$ for $u\ge1.4$ (consistent with fast decay); the
transition sits at $g_0\approx1/d$ as the trap picture requires, the
together and apart rows agree bin by bin (the $\iota$-symmetry made
visible), and summing bins reproduces $\psi\approx1.0/d$ with the
$\{g_0\lesssim1/d\}$ bins carrying essentially all of it.  Small-$u$
behaviour $1-\varphi(u)\approx u^{1/2\text{--}0.6}$ (not the naive
linear touch-count) --- unexplained, low stakes.

\emph{(c) The iid-ensemble constant flattens.}  Extending the
session-7 drift $4\E[\mathrm{bias}^2]\cdot k=0.95\to1.44$
($k=8..256$): $k=512$: $1.448\pm0.229$; $k=1024$: $1.171\pm0.324$
(discrete code, $20$M/$40$M trials) --- constant $\approx1.3$--$1.45$,
\emph{no} $\log k$; the budget question of session 7 (8d(b)) is
settled: no log factor.

\emph{(d) The one-point exponent is exactly $2$.}  From the clean
start (single curve, $x=y=\tfrac12$):
$f(m)=+0.006280(50),\ +0.002951(41),\ +0.001713(35),\
+0.000814(29),\ +0.000458(25)$ at $m=8,12,16,24,32$;
$f(m)\cdot m^2=0.402,0.425,0.438,0.469,0.469$ and the local exponent
reaches $2.00$ on $m=24\to32$: $f(m)\to0.47\,m^{-2}$.  This confirms
the trap-creation mechanism: a clean start must \emph{create} its
trap (Lemma~\ref{lem:recharge}(iii) prices a depth-$1/m$ trap at
$\asymp1/m^2$), whereas the stationary start carries a ready-made
linear atom --- hence $1/m^2$ against the stationary $1/d$; it also
resolves session 7's 8d(c) (the value was below MC resolution
there).  \emph{Session-9 correction:} the same data at odd $m$ shows
$f(m)$ \emph{alternates} in sign, so the leading $1/m^2$ term is the
alternating one and not the frozen-label trap term; see
Lemma~\ref{lem:parity}(a) and Theorem~\ref{thm:exactonepoint}.

%----------------------------------------------------------------------
\subsection{Session 9: the gap law is not needed; parity; an exact
one-point formula}\label{sec:s9}

\subsubsection{Short horizon: Problem~\ref{prob:gaplaw} is
unnecessary}\label{sec:shorthorizon}

The reduction of \S\ref{sec:upper} asked for a gap law uniform in $j$.
But in Theorem~\ref{thm:factor} the initial law $\mu_j$ enters only for
$j\le k$ while $d\ge k/8$, i.e.\ $j\le8d$: the chain has run for at
most $8d$ steps before the $d$ fresh steps begin, and over such a short
horizon the state-independent creation flux of
Lemma~\ref{lem:recharge}(iii) alone controls the gap.

\begin{proposition}[short-horizon gap bound]\label{prop:shorthorizon}
Let $g_0^{\mathrm{ini}}=\min(|u-v|,1-|u-v|)$ be the initial separation
($=\tfrac12$ in the antipodal ensemble) and let $g_j$ be the gap of the
marked state after $j$ iid chords.  Then for every $\varepsilon>0$ and
every $j\ge0$,
\[
\Pr_{\mu_j}\bigl[g_j\le\varepsilon\bigr]
\;\le\;\mathbf 1\{g_0^{\mathrm{ini}}\le\varepsilon\}\;+\;4\varepsilon^2 j .
\]
\end{proposition}

\begin{proof}
If $g_j\le\varepsilon<g_0^{\mathrm{ini}}$ there is a first step
$t\le j$ at which the gap enters $(0,\varepsilon]$ from above; by
Lemma~\ref{lem:recharge}(iii) each step does so with probability at
most $4\varepsilon^2$ whatever the current state (the four channels
--- shrink of either coordinate, either end of a flip --- each
contribute at most $\varepsilon^2$; merges only increase coordinates
and environment-only moves leave them unchanged).  Sum over $t\le j$.
\end{proof}

\emph{Status: PROVED.  VERIFIED (\texttt{s9\_survival flux}): at
$j=8$ the ratio $\Pr[g_j\le\varepsilon]/(4\varepsilon^2j)$ is
$1.00\pm0.05$ for $\varepsilon\le2^{-8}$ --- the bound is
\emph{sharp} at small $\varepsilon$ (each channel's entry probability
is exactly $\varepsilon^2$ while all coordinates exceed
$\varepsilon$, and nothing has had time to escape); at $j=64$ and
$j=512$ the ratio is $\le0.93$ resp.\ $\le0.92$ in every bin and
decreases to the stationary linear law $4\varepsilon$ at moderate
$\varepsilon$.}

\begin{corollary}[$\kappa$-free budget]\label{cor:kappafree}
Suppose Problem~\ref{prob:survival} holds with exponent $\alpha$:
$\E_{\mu_j}[\Phi_d^2;\,g_j\ge\varepsilon]\le C(\varepsilon d)^{-2\alpha}$
for all $j\le8d$.  Then, in the antipodal ensemble,
\[
\psi(j,d)\;\le\;32\,\varepsilon^2d+C(\varepsilon d)^{-2\alpha}
\;\le\;C'\,d^{-\alpha/(1+\alpha)}\qquad(j\le8d),
\]
i.e.\ Problem~\ref{prob:mixing} holds with
$c'=\alpha/(2(1+\alpha))$, and the budget $c'>1/11$ of
Corollary~\ref{cor:psienough} is met as soon as $\alpha>2/9$.  In
particular the expected value $\alpha=1/2$ gives $c'=1/6$ (against
$c'=1/4$ with a linear gap law), and Problem~\ref{prob:gaplaw} is
\emph{not needed} for the iid model.
\end{corollary}

\begin{proof}
Combine Proposition~\ref{prop:shorthorizon} ($j\le8d$) with
$|\Phi_d|\le1$ in the decomposition of \S\ref{sec:upper} and optimise
$\varepsilon=d^{-(1+2\alpha)/(2+2\alpha)}$; $c'>1/11$ iff
$11\alpha>2+2\alpha$.
\end{proof}

\emph{Status: PROVED.  The whole ``no-dust''/escape discussion of
\S\ref{sec:upper} is thereby demoted to a route for the \emph{sharp}
exponent $c=1/2$ (or for the stationary-start statement); the iid
model needs only Problem~\ref{prob:survival} with any $\alpha>2/9$.
In the true ensemble the indicator term is the short-arc atom already
accounted for as the $O(1/k)$ term of \S\ref{sec:annealed}.}

\begin{lemma}[$s$-trap flux]\label{lem:strap}
From any state, the marked total $s=a+b$ enters $(0,\sigma]$ from above
with probability at most $2\sigma^2$ per step; hence
$\Pr_{\mu_j}[s_j\le\sigma]\le2\sigma^2j$ in the antipodal ensemble.
\end{lemma}
\begin{proof}
Flips and merges do not decrease $s$; a shrink of $a$ lands at
$s'=a'+b\le\sigma$ only if $b\le\sigma$ and $a'\le\sigma-b$, which has
probability $\le(\sigma-b)^2\le\sigma^2$ (the shrink channel of
Lemma~\ref{lem:recharge}(iii)); likewise for $b$.
\end{proof}
\emph{Status: PROVED.  Consequently in Corollary~\ref{cor:kappafree}
one may also assume $s_0\ge\sigma$ at cost $16\sigma^2d$, so the
target estimate of Problem~\ref{prob:survival} need only be proved
for starts with $a_0,b_0\ge\varepsilon$ and $s_0\ge\sigma\asymp\varepsilon$.}

\subsubsection{Parity: the label is a function of time and the
environment block count}\label{sec:parity}

\begin{lemma}[block-count parity]\label{lem:parity}
Let $B_t$ be the total number of curves and $E_t$ the number of
environment curves (those through neither $u$ nor $v$) at time $t$.
Every step is a split or a merge, so $B_t\equiv B_0+t\pmod 2$; since
$B_t=E_t+1+\mathbf 1\{u\not\sim v\}$,
\[
\varepsilon_t\;=\;-(-1)^{B_0+t}\,(-1)^{E_t},
\qquad
(-1)^{N_t}\;=\;(-1)^{t}\,(-1)^{E_t-E_0},
\]
with $N_t$ the number of flip moves.  Hence
$\Phi_d(\Gamma)=(-1)^d\,\E_\Gamma[(-1)^{E_d-E_0}]$.
\end{lemma}

\begin{proof}
Splits add a curve, merges remove one; a flip (separate or join)
changes $B$ by one and leaves $E$ unchanged, every other move leaves
the label and changes $E$ by one.
\end{proof}

\emph{Status: PROVED (trivial, but structurally decisive).  Three
consequences.}

\emph{(a) The sign of the one-point function alternates.}  From the
clean start, $f(m)=\E[\varepsilon(F_m)]$ obeys
$f(m)\cdot(-1)^m>0$: measured $f(7)=-0.00753(22)$, $f(8)=+0.00599(22)$,
$f(9)=-0.00526(22)$, $f(11)=-0.00353$, $f(12)=+0.00295$,
$f(13)=-0.00275$, $f(15)=-0.00192$, $f(16)=+0.00213$,
$f(17)=-0.00149$ ($2\times10^7$ trials each).  Session 8 sampled only
even $m$ and read $f(m)\approx0.47/m^2$ as the frozen-label trap
term; that reading is \emph{incomplete}: the leading term is the
alternating one, $f(m)\approx(-1)^m\,c_1/m^2$ with $m^2|f(m)|$
increasing through $0.40,0.43,0.44,0.47,0.47$ ($m=8..32$) and
plausibly $c_1=\tfrac12$ (Theorem~\ref{thm:exactonepoint} below gives
exactly $\tfrac12$ for the $(u,v)$-averaged function), while the
non-alternating (trap) part, $\tfrac12(f(m{-}1)+f(m{+}1))+f(m)$ at
even $m$, is $\lesssim2\times10^{-4}$ at $m=8$.  Nothing downstream
is affected ($\psi=\E\Phi^2$ is sign-blind), but the smooth object is
$(-1)^d\Phi_d$, not $\Phi_d$, and Remark~\ref{rem:renewal}'s heuristic
``$\Phi_d\approx\Pr[\tau>d]$'' must be read with this sign.

\emph{(b) An obstruction for coupling proofs.}  Any two words of the
same length whose final environments have the same block-count parity
have the same flip parity; in particular two trajectories from the
same start that agree in marked geometry \emph{and} environment at
time $d$ have the same label.  A coupling proof of
Problem~\ref{prob:survival} must therefore produce, at time $d$,
environments of \emph{different} block-count parity with the same
law --- e.g.\ one carrying an extra tiny block.  The natural ``swap a
flip for a shrink--re-merge pair and let the environment take a tiny
extra split'' coupling has success probability $O(1/d)$ per step
and is repeated $d$ times, so it yields only $|\Phi_d|\le1-c$: no
decay.  Same-time $\iota$-mirroring is impossible ($\iota$ changes
$B$ by one); a one-step shifted mirror gives $\Phi_d\approx-\Phi_{d+1}$,
which is the alternation (a), not decay.  The parity of $N_d$ is thus
a genuinely ``continuous-state'' functional and the proof of
Problem~\ref{prob:survival} has to go through the renewal or spectral
route.

\emph{(c) First-flip tail from good starts is $1/t^2$, not $1/t$.}
From a stationary environment conditioned on $g_0\in[1/16,1/8)$
(\texttt{s9\_survival phic}), $t^2\Pr[\tau>t]\to11.1$ ($t=128,256$;
$16.1$ at $t=64$, $29$ at $t=32$); from the clean start with
$g_0=1/4$, $t\Pr[\tau>t]$ halves per doubling of $t$ from $t=32$ on.
The renewal heuristic explains this: a trap at depth $\eta\le\varepsilon$
is created at rate $8\eta\,d\eta$ (Lemma~\ref{lem:recharge}(iii)) and
survives $m$ further steps without flipping with probability
$\approx e^{-2\eta sm}$, so the kernel $K(m)=\int_0^\varepsilon
8\eta e^{-2\eta sm}d\eta\approx2/(s^2m^2)$ for $m\gg1/(\varepsilon s)$,
with total mass $\approx4\varepsilon/s<1$: the first-flip time is a
\emph{subcritical} renewal sum with a $1/t^2$ kernel, hence
$\Pr[\tau>t]\asymp t^{-2}$ and, if the alternating-renewal
cancellation holds, $|\Phi_d|\lesssim(\varepsilon d)^{-2}$ off small
gaps --- i.e.\ the true $\alpha$ in Problem~\ref{prob:survival} is
plausibly $2$, far above the $2/9$ needed.  Consistently,
$\E[\Phi_d^2\mid g_0=u/d]$ is at the noise floor ($\lesssim10^{-3}$,
often negative estimates) for every bin $u\ge3$ in the session-8
\texttt{psibar} logs at $d=128$--$1024$.  COMPUTED, not proved.

\subsubsection{An exact one-point formula}\label{sec:exactonept}

Averaging the marked points over the circle turns the one-point
function into a class function of the discrete precursor, and hook
characters compute it exactly.

\begin{theorem}[exact $(u,v)$-averaged one-point function]
\label{thm:exactonepoint}
Let $\sigma_m=\gamma\,\tau_1\cdots\tau_m$ with $\gamma$ a uniform
$n$-cycle and $\tau_i$ iid uniform transpositions, and let $c_i$ be
its cycle lengths.  Then
\[
\E\Bigl[\sum_i c_i(\sigma_m)^2\Bigr]
=\frac{n(n+1)}2+\sum_{k=1}^{n-1}(n-k)\Bigl(1-\frac{2k}{n-1}\Bigr)^{m},
\]
so that for uniform distinct $u,v$,
$\Pr[u\sim v\text{ in }\sigma_m]=\bigl(\E\sum c_i^2-n\bigr)/(n(n-1))$.
Consequently, for the continuum uniform split--merge chain
$(p_i(m))$ started from a single block,
\[
\E\Bigl[\sum_i p_i(m)^2\Bigr]
=\frac12+\frac14\Bigl[\frac{1+(-1)^m}{m+1}+\frac{1-(-1)^m}{m+2}\Bigr]
=\begin{cases}\tfrac12+\tfrac1{2(m+1)},& m\text{ even},\\[2pt]
\tfrac12+\tfrac1{2(m+2)},& m\text{ odd},\end{cases}
\]
and the $(u,v)$-averaged one-point function
$\bar f(m)=\E_{u,v}\E[\varepsilon(F_m)]=2\E\sum p_i^2-1$ equals
$1/(m+1)$ for even and $1/(m+2)$ for odd $m$:
$\bar f(m)=\dfrac{2m+3}{2(m+1)(m+2)}+(-1)^m\dfrac{1}{2(m+1)(m+2)}$.
\end{theorem}

\begin{proof}
Right multiplication by a uniform transposition acts on the
$\lambda$-isotypic component of $L^2(S_n)$ by
$r_\lambda=\chi_\lambda(\tau)/d_\lambda$, and the law of $\gamma$ is
the class measure of $n$-cycles, so for any $F$,
$\E F(\sigma_m)=\sum_\lambda \rho_\lambda\,r_\lambda^{m}\,
\langle\chi_\lambda,F\rangle$ with
$\rho_\lambda=\chi_\lambda(n\text{-cycle})$ and
$\langle\chi,F\rangle=\frac1{n!}\sum_\sigma\chi(\sigma)F(\sigma)$
(Fourier inversion for class measures).  Only hooks
$\lambda_k=(n-k,1^k)$ have $\rho_{\lambda}\ne0$, with
$\rho_{\lambda_k}=(-1)^k$, and their content sum gives
$r_{\lambda_k}=\bigl[(n-k)(n-k-1)-k(k+1)\bigr]/(n(n-1))=1-2k/(n-1)$.
For $F=\sum_ic_i^2$ use the hook generating function
$\sum_k\chi_{\lambda_k}(\mu)t^k=(1+t)^{-1}\prod_i(1-(-t)^{\mu_i})$ and
the exponential formula with one marked part:
\[
\sum_{n}x^n\sum_k t^k\langle\chi_{\lambda_k},F\rangle
=\frac1{1+t}\Bigl[\sum_{j\ge1}j\,(1-(-t)^j)x^j\Bigr]
\exp\Bigl(\sum_{j\ge1}\frac{(1-(-t)^j)x^j}{j}\Bigr)
=\frac1{1+t}\Bigl[\frac{x}{(1-x)^2}+\frac{tx}{(1+tx)^2}\Bigr]
\frac{1+tx}{1-x}.
\]
Extracting $[x^n]$: the first term gives
$\bigl[\binom{n+1}2+t\binom n2\bigr]/(1+t)$, the second
$t(1-(-t)^n)/(1+t)^2$; the coefficient of $t^k$ ($1\le k\le n-1$) is
$(-1)^kn+(-1)^{k-1}k=(-1)^k(n-k)$, and $\binom{n+1}2$ at $k=0$.
Multiply by $\rho_{\lambda_k}=(-1)^k$ and $r^m$.  The continuum limit
is the Riemann sum $\int_0^1(1-x)(1-2x)^m dx$, evaluated by
$y=1-2x$; the identification of $\sigma_m/n$ with the split--merge
chain from one block is Lemma~\ref{lem:splitmerge} (the $n$ points
carry arcs of length $1/n$).
\end{proof}

\emph{Status: PROVED.  VERIFIED: the finite-$n$ formula
(\texttt{s9\_hooks.py}, exact rationals; the inner products
$(-1)^k(n-k)$ were also checked against the full character sums for
$n\le12$) gives at $n=40$: $0.6007,\,0.5567,\,0.5314$ at $m=4,8,16$
against the continuum simulator's $0.6005,\,0.5561,\,0.5297$
(\texttt{s8\_splitmerge traj}), and the continuum values
$m=1:\ 2/3$, $m=2:\ 2/3$ agree with direct hand integration.
Remarks: (i) this is the $(u,v)$-averaged version of the one-point
function of \S\ref{sec:mixnum}(d); its non-alternating part
$\approx1/m$ is the frozen-label trap term fed by the linear initial
gap density (2 on $(0,\tfrac12)$) of a random separation, absent for
antipodal $u,v$, and its alternating part $(-1)^m/(2(m+1)(m+2))$ has
exactly the coefficient $\tfrac12$ towards which the antipodal
$m^2|f(m)|$ appears to converge; (ii) in spectral terms the averaged
start has spectral density $(1+y)/4$ on $y=1-2x\in[-1,1]$ --- linear
at \emph{both} ends $y=\pm1$, the two ends being the trap ($y=+1$)
and the alternating ($y=-1$, the $(-1)^{B}$ eigenfunction) edges;
(iii) the same technique cannot reach $\psi$ (a two-time quantity is
not a class function of the increment), but it does prove the
averaged one-point decay $\bar f(m)\le1/(m+1)$ outright, and gives
the exact law of $\E\sum p_i^2$ along the whole approach to
$\mathrm{PD}(1)$.}

\subsubsection{The spectral route on the quotient: twisted kernel}
\label{sec:twisted}

Lemma~\ref{lem:parity} and Corollary~\ref{cor:quotient} make the
spectral formulation of \S\ref{sec:spectral} concrete on the quotient
chain.  Write $P=P_{\mathrm{nf}}+P_{\mathrm f}$ for the non-flip and
flip parts of the quotient kernel and let
$P_\pi:=P_{\mathrm{nf}}-P_{\mathrm f}$ (the kernel twisted by
$(-1)^{F}$).  Then $\Phi_d=P_\pi^{\,d}\mathbf 1$, and since flips are
reversed by flips, $P_\pi$ is self-adjoint on $L^2(\tilde\pi)$
($\tilde\pi$ the quotient stationary law), so
\[
\psi_\pi(d)=\langle\mathbf 1,P_\pi^{2d}\mathbf 1\rangle
=\int_{[-1,1]}\lambda^{2d}\,\nu(d\lambda),\qquad
\nu=\text{spectral measure of }\mathbf 1\text{ under }P_\pi,\ \nu([-1,1])=1,
\]
and Problem~\ref{prob:mixing} at stationarity reads
$\nu(\{|\lambda|\ge1-\delta\})\le C\delta^{\beta}$
($\psi_\pi(d)\le Cd^{-\beta}$).  Two facts, one numerical and one a
proof sketch:

\emph{(a) No mass at the $-1$ end at stationarity (COMPUTED).}
$(-1)^{E}$ is an exact $\lambda=-1$ eigenfunction of $P_\pi$
(Lemma~\ref{lem:parity}), so the $2$-periodicity caveat of
\S\ref{sec:spectral} is a priori relevant.  But
$\langle P^d\varepsilon,P^{d+1}\varepsilon\rangle_\pi
=\int\lambda^{2d+1}d\nu$ versus $\int\lambda^{2d}d\nu$
(\texttt{s9\_survival oddeven}): $0.0559/0.0575=0.971$ at $d=16$ and
$0.0146/0.0150=0.974$ at $d=64$ (both $\pm0.03$), so the relative
mass of $\nu$ near $-1$ is $\lesssim3\%$ and compatible with zero; in
the continuum $E=\infty$ under $\tilde\pi$ and the eigenfunction
$(-1)^E$ is not in $L^2(\tilde\pi)$.  The sign alternation of
\S\ref{sec:parity}(a) is a finite-block (clean-start) phenomenon,
invisible at stationarity.

\emph{(b) A Cauchy--Schwarz/smoothing route to $\beta=1/2$
(SKETCH).}  The Dirichlet form of $P_\pi$ is
$\mathcal E(F)=\tfrac12\E[(F-F')^2;\mathrm{nf}]+\tfrac12\E[(F+F')^2;\mathrm f]$.
Let $k(x)=2ab$ be the flip rate and $Q(x,\cdot)=P_{\mathrm f}(x,\cdot)/k(x)$
the normalised flip kernel, which acts on each segment
$\{(a',s-a';\mathrm{env})\}$ separately by the trapezoidal density
$a'\mapsto|\{\alpha\in(0,a):a'-\alpha\in(0,s-a)\}|/(a(s-a))$.  By
Cauchy--Schwarz in the flip term,
$\mathcal E(F)\ge\tfrac12\int\tilde\pi(dx)\,k(x)\,(F+QF)^2(x)$, so
$\mathcal E(F)\le\delta\|F\|^2$ forces $F\approx-QF$ in
$L^2(\tilde\pi k)$, hence $F\approx Q^2F$: $F$ is nearly constant,
then nearly equal to minus itself, on every segment away from its
ends --- i.e.\ $\int_{\{ab\ge\eta\}}F^2d\tilde\pi\le C\delta\|F\|^2/\eta$
once the one-dimensional smoothing inequality for $Q^2$ on a segment
is quantified (the trapezoid density is $\ge c/s$ away from the
ends, $\asymp a'/s^2$ near them).  With $\tilde\pi(ab<\eta)\le C\eta$
(Proposition~\ref{prop:gaplawstat} plus the quadratic $s$-law),
$\langle F,\mathbf 1\rangle^2\le C\|F\|^2(\delta/\eta+\eta)$, and
$\eta=\sqrt\delta$ gives $\nu([1-\delta,1])\le C\sqrt\delta$, i.e.\
$\psi_\pi(d)\le Cd^{-1/2}$: $c=1/4>1/11$ \emph{at stationarity}.
Gaps: (i) the segment smoothing inequality (a weighted Poincar\'e
for $Q^2$ with weight $a'(s-a')$); (ii) the $-1$ end, to be handled
either by (a) made rigorous or by running the argument for
$\langle F,(I+P_\pi)F\rangle$; (iii) the passage from $\pi$ to
$\mu_j$, $j\le8d$ --- the short-horizon Proposition~\ref{prop:shorthorizon}
handles the gap coordinate but not the rest of the state, and
Lemma~\ref{lem:freeze} only trades $d$ for $j$.  (iii) is the real
obstacle for the iid model; a dust-insensitive comparison of
$\mu_{j}$ with $\pi$ for $j\ge d/2$ (a block of mass $w$ influences
$d$ steps with probability $\le2wd$) is the natural tool.

%----------------------------------------------------------------------
\subsection{Session 10: the block count is Harer--Zagier, the environment
is healthy on the short horizon, the exact size-biased law, and the
Markov-renewal form of the survival problem}\label{sec:s10}

Session 10 was to fix the exact statement behind
Problem~\ref{prob:survival} (TODO~10a, step~0) and choose a route.  The
first finding is negative and structural: \emph{the environment cannot
be avoided} --- without merges the parity of the flip count from a start
with $\min(a_0,b_0)=\varepsilon$ decays on the time scale
$\varepsilon^{-2}$, not $\varepsilon^{-1}$, and the budget of
Corollary~\ref{cor:kappafree} forces $\varepsilon\ll d^{-1/2}$, so a
proof must use the recharge of small marked coordinates by
environment blocks (\S\ref{sec:noenv}).  The second finding removes the
obstacle this creates.  The number of curves $B_t$ after $t$ chords has
an \emph{exact} law --- it is the Harer--Zagier vertex-count
distribution of a random gluing of a $2t$-gon --- with
$\E[3^{B_t}]=3+4t+2t^2$; a Chernoff bound then gives $B_t\le C\log d$
simultaneously for all $t\le9d$ except with probability $d^{-\gamma}$
(\S\ref{sec:blockcount}).  On that event every state on the horizon
has $\sum p_i^2\ge1/(C\log d)$ and an environment block of mass
$\ge(1-s)/(C\log d)$: the ``no-dust'' question is closed, up to
logarithms, in exactly the form the short horizon needs.  The same hook
computation gives the exact law of the mass of the curve through a
uniform point at every time (\S\ref{sec:sizebiased}; this contains
Theorem~\ref{thm:exactonepoint} and closes TODO~10b(b)).  Finally the
survival problem is put in Markov-renewal form
(\S\ref{sec:markovrenewal}): an exact operator identity for the
generating function $\sum_d\Phi_dz^d$ separates the two remaining
contents --- moments of the first-flip time (traps) and invertibility
of $I+K_z$ (cancellation) --- and a conditional decay proposition
states precisely what is left.

\subsubsection{The environment cannot be avoided}\label{sec:noenv}

Consider the marked pair alone, with merges suppressed (an adversarial
``dust'' environment): per step, $a\to a\sqrt U$ with probability
$a^2$, $b\to b\sqrt U$ with probability $b^2$, a flip
$(a,b)\to(\alpha+\beta,s-\alpha-\beta)$ with probability $2ab$, and
nothing otherwise.  Since merges only increase marked coordinates and
are the only moves that do, this is the chain that any argument using
the environment solely through upper bounds would have to control.
Its scaling is elementary: all rates are quadratic in the coordinates,
so from $(a_0,b_0)=(\varepsilon,\varepsilon)$ nothing happens for
$\asymp\varepsilon^{-2}$ steps, and from $(c_0,\varepsilon)$ the large
coordinate decays like $t^{-1/2}$ before a flip (rate $2\varepsilon b$)
can occur, the first flip arriving only at $t\asymp\varepsilon^{-2}$.
Measured (\texttt{s10\_nomerge}, $\Phi_t=\E[(-1)^{N_t}]$):
from $(\tfrac14,\tfrac14)$, $\Phi_t=0.130,\,0.031,\,0.0070,\,0.0024$
at $t=8,16,32,64$ (then noise floor $10^{-3}$; $t^2\Phi_t\approx8$--$10$
on $16\le t\le64$) --- a healthy $1/t^2$-like decay without any
environment; but from $(\tfrac14,0.02)$,
$\Phi_t=0.53,\,0.25,\,0.088,\,0.040$ at $t=64,256,1024,4096$, and from
$(0.02,0.02)$, $\Phi_t=0.90,\,0.68,\,0.30,\,0.13$ at
$t=64,256,1024,8192$: decay only on the scale $\varepsilon^2t$, and
slow even there.  In the budget of Corollary~\ref{cor:kappafree} the
trap-entry flux forces $\varepsilon\ll d^{-1/2}$, hence
$\varepsilon^2d\ll1$: a bound depending on the environment only
through upper bounds is worthless.  \emph{Status: COMPUTED (the scaling
statement is trivial).  Consequence: any proof of
Problem~\ref{prob:survival} must lower-bound the masses of environment
blocks over the horizon --- the recharge mechanism of
Remark~\ref{rem:renewal}.  The next subsection supplies exactly this.}

\subsubsection{The exact law of the number of curves}\label{sec:blockcount}

Let $B_t$ be the number of curves after $t$ iid chords (the total
block count of the marked chain from the one-block start; it does not
depend on the marks).

\begin{theorem}[the block count is Harer--Zagier]\label{thm:blockcount}
For every $x$,
\[
\sum_{t\ge0}\E\bigl[x^{B_t}\bigr]z^t
=\frac1{2z}\Bigl[\Bigl(\frac{1+z}{1-z}\Bigr)^{x}-1\Bigr],
\qquad\text{i.e.}\qquad
\Pr[B_t=m]=[z^{t+1}]\,\frac{(2\operatorname{artanh}z)^m}{2\,m!},
\]
so $B_t\equiv t+1\pmod2$, $\Pr[B_t=t+1-2g]=\varepsilon_g(t)/(2t-1)!!$
with $\varepsilon_g(t)$ the Harer--Zagier numbers \cite{HZ86}, and for
integer $x\ge1$ the expectation $\E[x^{B_t}]$ is a polynomial in $t$ of
degree $x-1$:
\[
\E[2^{B_t}]=2+2t,\qquad \E[3^{B_t}]=3+4t+2t^2,\qquad
\E[4^{B_t}]=4+\tfrac{20}3t+4t^2+\tfrac43t^3=\tfrac43(t+1)(t^2+2t+3),\qquad
\Pr[B_t=1]=\frac{1+(-1)^t}{2(t+1)}.
\]
\end{theorem}

\begin{proof}
The $2t$ endpoints are iid uniform on the circle, a.s.\ distinct; by
exchangeability, conditionally on their positions the pairing into
chords is uniform among the $(2t-1)!!$ pairings.  Refining $\gamma$ at
the endpoints does not change the curve partition
(Lemma~\ref{lem:splitmerge}), so $B_t$ is the number of cycles of
$\gamma_{2t}\tau$, with $\gamma_{2t}$ the rotation of the $2t$
endpoints and $\tau$ the pairing involution --- the number of vertices
of the surface obtained by gluing the sides of a $2t$-gon along a
uniform pairing.  Harer--Zagier's formula \cite{HZ86} is
$1+2\sum_{t\ge0}\bigl[\sum_g\varepsilon_g(t)N^{t+1-2g}\bigr]
z^{t+1}/(2t-1)!!=\bigl((1+z)/(1-z)\bigr)^N$; dividing by $(2t-1)!!$ turns
the bracket into $\E[N^{B_t}]$, which is the displayed identity for
$x=N$, and both sides are polynomials in $N$ for each $t$.  Expanding
$((1+z)/(1-z))^x=\exp(2x\operatorname{artanh}z)$ in powers of $x$ gives
the law; the polynomials are the coefficients of
$(1-z)^{-x}\sum_k\binom{x}{2k+1}z^{2k}$.
\end{proof}

\emph{Status: PROVED (the statement is classical once the chord
diagram is identified; the identification is the content).  A second
proof runs through the finite-$n$ hook machinery of
Theorem~\ref{thm:exactonepoint}: $x^{c(\sigma)}=\sum_\lambda
s_\lambda(1^x)\chi_\lambda(\sigma)$, so
$\E[x^{c(\sigma_t)}]=\sum_{k}(-1)^k r_k^{\,t}\,s_{(n-k,1^k)}(1^x)$ with
$s_{(n-k,1^k)}(1^x)=x(x{+}1)\cdots(x{+}n{-}k{-}1)\cdot(x{-}1)\cdots(x{-}k)
/\bigl(n\,(n{-}k{-}1)!\,k!\bigr)$ by hook-content; for integer $x$ only
$k\le x-1$ survive, e.g.\ $\E[2^{c(\sigma_t)}]=(n+1)-(n-1)r_1^{\,t}$,
$\E[3^{c(\sigma_t)}]=\tfrac12(n+1)(n+2)-(n^2-1)r_1^{\,t}+\tfrac12(n-1)(n-2)r_2^{\,t}$,
whose $n\to\infty$ limits are $2+2t$ and $3+4t+2t^2$
(\texttt{s10\_blocks.py}, exact rationals, Lagrange-extrapolated in
$1/(n-1)$, gives the polynomials for $x\le5$).  VERIFIED: the full law
against a direct simulation of the split--merge chain
(\texttt{s10\_blocksim.py}, $2\times10^5$ trials, $t\le6$): every
probability agrees to $\pm0.001$, e.g.\ $\Pr[B_6=1,3,5,7]=
0.1429,0.6222,0.2222,0.0127$ predicted vs.\ $0.1431,0.6230,0.2212,0.0127$
measured.}

\begin{corollary}[block count over the horizon]\label{cor:horizon}
For all $T\ge1$ and $K\ge1$,
\[
\Pr\Bigl[\max_{0\le t\le T}B_t\ge K\Bigr]\;\le\;3^{\,1-K}(T+1)^3 .
\]
In particular, for $\gamma>0$ and
$K_\gamma(T):=\lceil(3+\gamma)\log_3(T+1)\rceil+1$, the event
$H_T=\{B_t<K_\gamma(T)\ \forall t\le T\}$ has probability
$\ge1-(T+1)^{-\gamma}$, and on $H_T$, for every $t\le T$: the total number of
blocks is $<K_\gamma(T)$, $\sum_ip_i(t)^2\ge1/K_\gamma(T)$, the
environment has fewer than $K_\gamma(T)$ blocks and hence a block of
mass $\ge(1-s_t)/K_\gamma(T)$.
\end{corollary}

\begin{proof}
Markov with $x=3$: $\Pr[B_t\ge K]\le3^{-K}(3+4t+2t^2)\le3^{1-K}(t+1)^2$;
sum over $t\le T$.  With $K=K_\gamma(T)$,
$3^{1-K}\le(T+1)^{-3-\gamma}$.  The consequences are Cauchy--Schwarz
($\sum p_i^2\ge1/\#\text{blocks}$) and pigeonhole.
\end{proof}

\emph{Status: PROVED.  In Problem~\ref{prob:survival} the start
$\Gamma_0\sim\mu_j$ ($j\le8d$) followed by $d$ steps is the chain from
one block over at most $9d$ steps, so $H_{9d}$ has probability
$\ge1-(9d+1)^{-\gamma}$ for any fixed $\gamma$ at the price
$K=O_\gamma(\log d)$.  This is the ``no-dust'' estimate that
\S\ref{sec:upper} and the session-8/9 handoffs listed as the open
ingredient of every route: over the short horizon it holds for free,
up to a logarithm, because the chain simply has not had time to create
many blocks --- $\E[B_t]=[z^t]\,(1+z)\operatorname{artanh}(z)/(z(1-z))=\log t+O(1)$
and the pgf makes the tail geometric.  (At stationarity the count is infinite and the analogous
statement is false as stated; the short horizon of
Theorem~\ref{thm:factor} is again what makes the iid model tractable.)
Since the required exponent in Corollary~\ref{cor:kappafree} is only
$\alpha>2/9$, losing logarithms is harmless.}

\subsubsection{The exact law of the block of a uniform point}
\label{sec:sizebiased}

\begin{theorem}[size-biased block mass at all times]\label{thm:sizebiased}
Let $U$ be uniform on the circle, independent of the chords, and let
$\zeta_m$ be the mass of the curve through $U$ after $m$ chords (chain
from one block).  Then
\[
\Pr[\zeta_m\in dx]=\bigl(1-(1-2x)^m\bigr)\,dx\ \text{ on }(0,1),
\qquad
\Pr[\zeta_m=1]=\frac{1+(-1)^m}{2(m+1)} ,
\]
and consequently, for every $r\ge1$,
\[
\E\Bigl[\sum_ip_i(m)^r\Bigr]=\E[\zeta_m^{\,r-1}]
=\frac1r+\int_0^1(1-x^{r-1})(1-2x)^m\,dx .
\]
By rotation invariance the same law governs the mass of the curve
through any fixed point, in particular through $u$ (and through $v$).
At finite $n$: $\Pr[|C_{\sigma_m}(1)|=j]=(1-r_j^{\,m})/n$ for $j<n$ and
$\Pr[|C_{\sigma_m}(1)|=n]=\frac1n\sum_{k=0}^{n-1}r_k^{\,m}$, with
$r_k=1-2k/(n-1)$; equivalently
$(-1)^k\langle\chi_{(n-k,1^k)},\sum_ic_i^{\,r}\rangle=n^{r-1}-k^{r-1}$
for $1\le k\le n-1$ and $\sum_{j\le n}j^{r-1}$ for $k=0$.
\end{theorem}

\begin{proof}
For $F$ on $\{1,\dots,n\}$ put $G_F(\sigma)=\sum_{u}F(|C_\sigma(u)|)
=\sum_ic_iF(c_i)$.  The exponential formula with one marked cycle, as
in the proof of Theorem~\ref{thm:exactonepoint} (weight $jF(j)$ for a
marked cycle of length $j$ replaces $j\cdot j$ there), gives
\[
\sum_nx^n\sum_kt^k\langle\chi_{(n-k,1^k)},G_F\rangle
=\frac1{1+t}\Bigl[\sum_{j\ge1}F(j)\bigl(1-(-t)^j\bigr)x^j\Bigr]\frac{1+tx}{1-x}.
\]
Extracting $[x^n]$: $\sum_{j<n}F(j)(1-(-t)^j)(1+t)+F(n)(1-(-t)^n)$;
dividing by $1+t$ and extracting $[t^k]$ gives, for $1\le k\le n-1$,
$-(-1)^kF(k)+(-1)^kF(n)$, and $\sum_{j\le n}F(j)$ for $k=0$.  Hence,
by the Fourier inversion of Theorem~\ref{thm:exactonepoint},
$\E[G_F(\sigma_m)]=\sum_{j\le n}F(j)+\sum_{k=1}^{n-1}r_k^{\,m}(F(n)-F(k))
=\sum_{j<n}F(j)(1-r_j^{\,m})+F(n)\sum_{k=0}^{n-1}r_k^{\,m}$.  Since
$G_F(\sigma)/n=\E_U[F(|C_\sigma(U)|)]$ for a uniform point $U$, this is
the finite-$n$ law; $F(j)=j^{r-1}$ gives the hook inner products, and
$j=nx$, $r_j^{\,m}\to(1-2x)^m$ gives the continuum statements
(Lemma~\ref{lem:splitmerge} for the identification).
\end{proof}

\emph{Status: PROVED.  VERIFIED: the hook identity
$(-1)^k\langle\chi_k,\sum c^r\rangle=n^{r-1}-k^{r-1}$ by brute-force
character sums for $r\le6$, $n\le9$ (\texttt{s10\_hooks34.py}); the
continuum law and the moments $r=2,3,4$ against simulation
(\texttt{s10\_sizebiased.py}, $2\times10^5$ trials): at $m=8$ the ten
decile masses agree to $\pm0.001$, the atom $0.1115$ vs $1/9$, and
$\E\sum p^{2,3,4}=0.5561,0.3947,0.3140$ vs $0.5556,0.3939,0.3131$.
Remarks.  (i) $r=2$ is Theorem~\ref{thm:exactonepoint}; the atom is
$\Pr[B_m=1]$ of Theorem~\ref{thm:blockcount}, as it must be.  (ii) The
drift identity $\E[\Delta L]=2L^2-\tfrac73\sum p_i^4$ ($L=\sum p_i^2$)
is now closed in expectation along the whole approach to
$\mathrm{PD}(1)$: $\E\sum p_i^4(m)=\tfrac14+\int_0^1(1-x^3)(1-2x)^mdx$.
(iii) Small marked blocks: $\Pr[\text{mass of }u\text{'s curve}\le\eta]
=\int_0^\eta(1-(1-2x)^m)dx\le\min(\eta,\,m\eta^2)$ --- the
$m\eta^2$ regime is the creation flux of Lemma~\ref{lem:recharge}(iii)
made exact for the block through a fixed point, and the crossover at
$\eta\asymp1/m$ is the trap scale.  (iv) The density $1-(1-2x)^m$ is
$\ge1-e^{-2mx}$ on $(0,\tfrac12]$: the size-biased pick is macroscopic with probability
bounded below uniformly in $m\ge1$ --- the quantitative handle for
recharge that TODO~10a(A)(iii) asked for, now superseded by
Corollary~\ref{cor:horizon}.}

\subsubsection{Recharge under the block bound}\label{sec:recharge}

Fix $K\ge3$ and let $H=H_T$ be the event of Corollary~\ref{cor:horizon}
with $K_\gamma(T)$ replaced by $K$.

\begin{lemma}[recharge of the marked total]\label{lem:srecharge}
Let $K\ge3$, $s_0\ge\varepsilon$ with $\varepsilon\le1/(2K)$.  Then for every
$T\ge1$,
\[
\Pr\Bigl[\max_{t\le T}s_t<\tfrac1{2K}\Bigr]
\;\le\;e^{-\varepsilon T/K}+2\varepsilon^2T+\Pr[H^c].
\]
The same holds for each marked coordinate separately: if
$a_0\ge\varepsilon$ then
$\Pr[\max_{t\le T}a_t<1/(2K)]\le e^{-\varepsilon T/(3K)}+2\varepsilon^2T+\Pr[H^c]$,
and likewise for $b$.
\end{lemma}

\begin{proof}
On $H\cap\{s_t<1/(2K)\}$ the environment has mass $>1-1/(2K)\ge\tfrac12$
in fewer than $K$ blocks, so it has a block of mass $z\ge1/(2K)$; the
chord with one cut in the marked material and one in that block has
probability $2s_tz\ge\varepsilon/K$ as long as $s_t\ge\varepsilon$, and
it raises $s$ by $z\ge1/(2K)$.  Flips and merges never decrease $s$, and
$s$ enters $(0,\varepsilon]$ from above with probability $\le2\varepsilon^2$
per step (Lemma~\ref{lem:strap}); so the event $\{\max_{t\le T}s_t<1/(2K)\}$
is contained in $\{s_t<\varepsilon\text{ for some }t\le T\}\cup H^c\cup
\{T\text{ consecutive failures of an event of probability }\ge\varepsilon/K\}$.
For a single coordinate $a$: if $s_t\le\tfrac12$ the same merge
argument applies to $a$ (probability $\ge2a\cdot1/(2K)\ge\varepsilon/K$);
if $s_t>\tfrac12$ and $a_t<1/(2K)$ then $b_t\ge\tfrac13$, a flip has
probability $2a_tb_t\ge2\varepsilon/3$ and lands at
$a'=\alpha+\beta\ge\beta\sim U(0,b_t)$, so $a'\ge1/(2K)$ with conditional
probability $\ge1-3/(2K)\ge\tfrac12$: success probability
$\ge\varepsilon/3\ge\varepsilon/(3K)$ per step in both regimes; the
entry flux into $(0,\varepsilon]$ is $\le2\varepsilon^2$ per step
(shrink of $a$, flip landing small; Lemma~\ref{lem:recharge}(iii)).
\end{proof}

\emph{Status: PROVED.  With $T=(3K/\varepsilon)\log(1/\delta)$ the
failure probability is $\delta+6\varepsilon K\log(1/\delta)+\Pr[H^c]$,
small precisely when $\varepsilon\ll1/K$ --- the regime in which
recharge is needed at all (for $\varepsilon\ge1/(2K)$ the start is
already log-macroscopic).  At the optimising
$\varepsilon=d^{-(1+2\alpha)/(2+2\alpha)}\in[d^{-3/4},d^{-1/2})$ of
Corollary~\ref{cor:kappafree} ($\alpha<1$) and $K=O(\log d)$,
$T=O(\varepsilon^{-1}\log^2d)\ll d$ and the failure probability is
$O(d^{-\gamma})+O(\varepsilon\log^2d)$, which is
$o\bigl((\varepsilon d)^{-2\alpha}\bigr)=o(d^{-\alpha/(1+\alpha)})$ since
$\varepsilon^{1+2\alpha}d^{2\alpha}=d^{-1/(2+2\alpha)}$.}

\begin{lemma}[log-drift of the marked total]\label{lem:logdrift}
On $H$, for a state with $s\le\tfrac12$,
\[
\E[\log s_{t+1}-\log s_t\mid\Gamma_t]\;\ge\;\frac{s}{K}\log\Bigl(1+\frac1{2Ks}\Bigr)-\frac{s^2}2 ,
\]
which is $\ge s\,(\log2-\tfrac14)/K>0$ for $s\le1/(2K)$; the negative
jumps of $\log s$ are those of $\tfrac12\log U$ (exponential tails),
the positive ones are unbounded above only through merges.
\end{lemma}

\begin{proof}
Merges: with probability $\ge2sz$ the total gains $z$, and the largest
environment block has $z\ge(1-s)/K\ge1/(2K)$; $\log(1+z/s)$ is
increasing in $z$.  Shrinks: $a\to a\sqrt U$ with probability $a^2$
changes $\log s$ by $\log(1-(a/s)(1-\sqrt U))\ge\log\sqrt U$, of mean
$-\tfrac12$; likewise $b$; $a^2+b^2\le s^2$.  Flips preserve $s$;
environment-only moves do not touch it.
\end{proof}

\emph{Status: PROVED (a calculation).  It says that, once the block
count is bounded, the marked total is a one-dimensional process with
positive log-drift below the level $\asymp1/K$ and light-tailed
down-jumps; standard drift arguments (Lyapunov function
$s^{-\theta}$, small $\theta$) then bound the occupation time of
$\{s<c/K\}$ over any window by a small fraction with polynomially
small failure probability.  This is the ``marked pair health''
input for the tail of the first-flip time in
\S\ref{sec:markovrenewal}; it is NOT written out.}

\subsubsection{The Markov-renewal form of Problem~\ref{prob:survival}}
\label{sec:markovrenewal}

Let $\tau$ be the first flip time of the quotient chain and, for
$|z|\le1$ and bounded measurable $F$,
\[
(K_zF)(x):=\E_x\bigl[z^{\tau}F(\Gamma_\tau)\bigr],
\qquad \|K_z\|_{L^\infty\to L^\infty}\le\sup_x\E_x|z|^{\tau}\le|z| ,
\]
since $\tau\ge1$.  $K_1$ is the post-flip kernel (Markov whenever
$\tau<\infty$ a.s.).

\begin{proposition}[generating function of the parity]\label{prop:genfun}
For $|z|<1$, as bounded functions of the initial state,
\[
\sum_{d\ge0}\Phi_d\,z^d\;=\;\frac1{1-z}\,(I+K_z)^{-1}(I-K_z)\mathbf 1 .
\]
\end{proposition}

\begin{proof}
With flip times $\tau_1<\tau_2<\cdots$, $\Phi_d=\sum_{k\ge0}(-1)^k
\Pr_x[\tau_k\le d<\tau_{k+1}]$ and
$\sum_dz^d\mathbf 1\{\tau_k\le d<\tau_{k+1}\}=(z^{\tau_k}-z^{\tau_{k+1}})/(1-z)$;
by the strong Markov property at the flip times
$\E_x[z^{\tau_k}]=(K_z^k\mathbf 1)(x)$.  Sum the Neumann series, which
converges since $\|K_z\|\le|z|<1$.
\end{proof}

\emph{Status: PROVED.  It is exact and separates the problem: the
factor $(I-K_z)\mathbf 1=\E_\cdot[1-z^{\tau}]$ carries the first-flip
tail (traps), the resolvent $(I+K_z)^{-1}$ carries the cancellation
(the parity of the number of flips per unit time).  Note
$(I-K_1)\mathbf 1=0$ wherever $\tau<\infty$ a.s.\ (on $\{ab>0\}$, since
$\sum_t2a_tb_t=\infty$ a.s.\ there --- taken for granted below), so the
right side is a difference quotient at $z=1$: the decay of $\Phi_d$ is the regularity of
$G(z):=(I+K_z)^{-1}(I-K_z)\mathbf 1$ on the closed disc.}

\begin{proposition}[conditional decay]\label{prop:condecay}
Let $\mathcal B$ be a Banach space of functions on the state space
containing $\mathbf 1$, on which the evaluations $F\mapsto F(x)$,
$x\in S'$, are uniformly bounded, and let $\gamma\in(0,1)$.  Suppose
\begin{enumerate}[leftmargin=2em]
\item[(R1)] $z\mapsto K_z\in\mathcal L(\mathcal B)$ is $C^{1,\gamma}$
on $\{|z|\le1\}$ (for a weighted sup-norm space
$\mathcal B=\{F:\sup|F|/V<\infty\}$ it suffices that
$\sup_x\E_x[\tau^{1+\gamma}V(\Gamma_\tau)]/V(x)<\infty$, by
$|\tau z_1^{\tau-1}-\tau z_2^{\tau-1}|\le2\tau^{1+\gamma}|z_1-z_2|^\gamma$);
\item[(R2)] $I+K_z$ is invertible on $\mathcal B$ for every $|z|\le1$,
with $\sup_{|z|\le1}\|(I+K_z)^{-1}\|<\infty$.
\end{enumerate}
Then $\sup_{x\in S'}|\Phi_d(x)|\le C\,d^{-\gamma}$.
\end{proposition}

\begin{proof}
Under (R1)--(R2), $G$ is $C^{1,\gamma}$ on the closed disc with
$G(1)=0$, so $\hat\Phi(z)=(G(z)-G(1))/(1-z)$ extends to a
$C^{0,\gamma}$ function on the closed disc; it is analytic inside, so
its boundary Fourier coefficients are the $\Phi_d$, and the Fourier
coefficients of a H\"older-$\gamma$ function are $O(d^{-\gamma})$.
Evaluate at $x\in S'$.
\end{proof}

\emph{\textbf{SESSION-11 CORRECTION.}  Condition (R2) is FALSE at
$z=\pm1$ for every function space containing the bounded functions:
$K_1\varepsilon=-\varepsilon$, because exactly one flip has occurred at
time $\tau$ (Proposition~\ref{prop:r2false}).  Proposition~\ref{prop:condecay}
is therefore vacuous as stated.  The same holds in the
label-free bookkeeping, with $\kappa_{-1}[(-1)^E]=-(-1)^E$.  It is
repaired in \S\ref{sec:s11}: (R2) is replaced by the two conditions
(C1) and (C2) of \S\ref{sec:s11left}, both statements about the single
Markov kernel $\kappa_1$, and everything on the closed disc away from
$z=\pm1$ is proved unconditionally by
Theorem~\ref{thm:idleloop}.  The discussion of (R1) below is
unaffected, and so is the reduction paragraph at the end of this
subsection.  Read the rest of this paragraph with that substitution.}

\emph{Status: PROVED as a reduction; (R1) and (R2) are OPEN.  What each
needs, and what is in hand:}

\emph{(R1) --- traps.  A uniform bound $\sup_{x\in S'}\Pr_x[\tau>t]\le
C_Kt^{-2}$ (measured: $t^2\Pr[\tau>t]\to11$ from stationary
environments, \S\ref{sec:parity}(c)) gives every moment $\tau^{1+\gamma}$,
$\gamma<1$, hence $\gamma$ arbitrarily close to $1$ in
Proposition~\ref{prop:condecay}, far above the $4/9$ that
Corollary~\ref{cor:kappafree} needs.  Ingredients: the state-independent
trap-creation flux (Lemma~\ref{lem:recharge}(iii)); the recharge of
Lemma~\ref{lem:srecharge} and the drift of Lemma~\ref{lem:logdrift},
both available on $H$ by Corollary~\ref{cor:horizon}; and the fact
that a trap of depth $g$ is left by a flip alone, at rate
$2g(s-g)$, needing no environment.  The proof of the $1/t^2$ tail is
a renewal computation with these inputs (the subcritical kernel
$\approx2/(s^2m^2)$ of \S\ref{sec:parity}(c)), which we have not
written.}

\emph{(R2) --- cancellation.  This is the parity content and cannot be
had from trap accounting.  It is needed on the whole closed disc: the
trivial bound $\|K_z\|\le|z|$ gives nothing on $|z|=1$.  At $z=1$ it
is the statement that $-1$ is not in the spectrum of the post-flip
chain $K_1$ on $\mathcal B$, with a uniform inverse; near $z=1$ it
follows by perturbation under (R1).  At other points of the circle it
is the non-lattice condition of Markov renewal theory --- no $F\ne0$
with $\E_x[z^{\tau}F(\Gamma_\tau)]=-F(x)$ --- and at $z=-1$, where
$K_{-1}F=\E_\cdot[(-1)^{\tau}F(\Gamma_\tau)]$, it says the parity of the
inter-flip times is not asymptotically deterministic; this is where
the $(-1)^d$ alternation of \S\ref{sec:parity}(a) lives and it needs
its own statement.  The Doeblin structure is on the marked pair: two consecutive
flips from any $(a,b)$ with total $s$ produce $a''$ whose law dominates
$\tfrac14\cdot\mathrm{Unif}[s/4,3s/4]$ (the trapezoid of
Lemma~\ref{lem:recharge}(ii) is symmetric and unimodal, so the first
flip lands in $[s/4,3s/4]$ with probability $\ge\tfrac12$, and from
there the second has density $\ge1/s$ on that interval).  This is a
minorisation on the fibre $\{a'+b'=s,\ \mathrm{env}\}$ only; the base
$(s,\mathrm{env})$ is not regenerated, so $\mathcal B$ must be a
weighted space in which the post-flip chain is $V$-geometrically
ergodic in the base coordinates as well (Meyn--Tweedie), with the
weight matched to the moment condition in (R1).  On $H$ the base is
finite-dimensional ($<K$ blocks), which is what makes this plausible.
Because the label parity is determined by the environment block-count
parity (Lemma~\ref{lem:parity}), no argument can avoid the environment
in (R2) either; but the environment enters only through the ergodicity
of the base chain between flips, not through any fine property of
$\mathrm{PD}(1)$.}

\emph{Reduction to Problem~\ref{prob:survival}.  Take
$S'=S'_K:=\{a,b\ge1/(16K)\}$ and enforce $H_{9d}$ (so $C=C(K)$ in
Proposition~\ref{prop:condecay} may depend on the block bound, which
holds up to time $d$).  Granting Proposition~\ref{prop:condecay} on
$S'_K$ with $C(K)$ polynomial in $K$, and granting the occupation-time
consequence of Lemma~\ref{lem:logdrift} (that after $s$ first reaches
$1/(2K)$ it stays $\ge1/(4K)$ for at least half of the next $T_2$
steps, except with polynomially small probability; NOT written out),
the chain from $\Gamma_0\sim\mu_j$ with $a_0,b_0\ge\varepsilon$ enters
$S'_K$ at a stopping time $T'$: first $s$ reaches $1/(2K)$ within
$T_1=O(\varepsilon^{-1}K\log d)$ steps (Lemma~\ref{lem:srecharge});
then, on the times with $s_t\ge1/(4K)$ and $\min(a_t,b_t)\ge\varepsilon$
(the latter failing with probability $\le4\varepsilon^2(T_1+T_2)$ by
the flux bound), a flip occurs with probability $\ge2\varepsilon(s_t-\varepsilon)
\ge\varepsilon/(4K)$ per step and lands with $\min(a',b')\ge s_t/4\ge1/(16K)$
with conditional probability $\ge\tfrac12$ (the trapezoid is
symmetric and unimodal), so $T'\le T_1+T_2$ with
$T_2=O(\varepsilon^{-1}K\log d)$ except with probability
$O(d^{-\gamma})+O(\varepsilon K\log d)$.  Note that a single
application of Lemma~\ref{lem:srecharge} to each coordinate does not
suffice: it makes $a$ and $b$ large at different times, and a
coordinate of size $1/(2K)$ shrinks on the scale $K^2\ll\varepsilon^{-1}K$.
By the strong Markov property
$|\Phi_d(\Gamma_0)|\le\E|\Phi_{d-T'}(\Gamma_{T'})|+\Pr[T'>T_1+T_2]
\le C(K)(d/2)^{-\gamma}+O(d^{-\gamma})+O(\varepsilon K\log d)$, and
with $K=O(\log d)$ and $\varepsilon$ as in Corollary~\ref{cor:kappafree}
this is $\le C(\varepsilon d)^{-2\alpha}$ for $\alpha<\gamma/2$ (for
$\varepsilon\ge1/(16K)$ no recharge is needed).  Thus
\emph{Problem~\ref{prob:survival}, hence the iid model, is reduced to
(R1) and (R2) on the log-macroscopic set $S'_K$ plus the
occupation-time step}, with the environment controlled throughout by
Corollary~\ref{cor:horizon}.}

%----------------------------------------------------------------------
\subsection{Session 11: (R2) is false; the idle loop proves its repair
on the whole disc except two points; and the residual condition is one
scalar statement about the parity of the environment block
count}\label{sec:s11}

Session 11 was directed at condition (R2) of
Proposition~\ref{prop:condecay} --- uniform invertibility of $I+K_z$ on
the closed unit disc --- under the instruction that a fourth
reformulation of the cancellation core would not count as progress.
There are four findings, in logical order.

\begin{enumerate}[leftmargin=2.2em]
\item[(1)] \textbf{The label is a passenger.}  The quotient chain
$\bar\Gamma=(p_1,p_2;\mathrm{env})$ --- the two marked masses and the
environment, with the label $\varepsilon$ forgotten --- is
\emph{autonomous}: its transition law does not depend on
$\varepsilon$, and the involution $J$ that swaps the two sectors
keeping $\bar\Gamma$ fixed commutes with the chain
(Lemma~\ref{lem:autonomy}).  Hence $K_z=\kappa_z\oplus(-\kappa_z)$ for
the quotient first-flip kernel $\kappa_z$, and
$\spec K_z=\pm\spec\kappa_z$.  One line of this gives
$\Pr_\pi[u\sim v]=\tfrac12$.
\item[(2)] \textbf{(R2) is false.}  With the label kept,
$K_1\varepsilon=-\varepsilon$ --- exactly one flip has occurred at
time $\tau$ --- so $I+K_1$ has a bounded null vector.  On the quotient,
where that artefact is absent, $I+\kappa_{-1}$ still has the bounded
null vector $(-1)^{E}$ ($E=$ number of environment curves), by
Lemma~\ref{lem:parity}.  Either way
$\sup_{|z|\le1}\|(I+K_z)^{-1}\|=\infty$ and
Proposition~\ref{prop:condecay} is vacuous as stated
(Proposition~\ref{prop:r2false}).
\item[(3)] \textbf{An unconditional contraction off $z=\pm1$.}  The
chain can waste exactly two steps by splitting a curve and merging the
two pieces it just created, returning to the same state without
flipping.  This ``idle loop'' has probability
$\lambda(\Gamma)=\tfrac13\sum_jr_j^4$ --- an exact formula, $r_j$ the
arc-refined mass list --- and it forces
$\|\kappa_z\|_\infty\le(1-\lambda)/|1-\lambda z^2|
\le1/(1+\lambda(1-\Re z^2))$ for all $|z|\le1$
(Theorem~\ref{thm:idleloop}).  On the block-bounded state space this is
uniform, with $\lambda\ge\tfrac13K^{-3}$.  The non-lattice condition of
Markov renewal theory on the \emph{whole} circle --- Pitfall~45,
TODO~11b(v), and the third name of the cancellation core --- is
therefore \emph{proved}, unconditionally and with explicit constants.
\item[(4)] \textbf{What is left is two scalar conditions on one Markov
kernel.}  After (3), $\widehat\Phi(z)=(1-z)^{-1}(I+\kappa_z)^{-1}
(I-\kappa_z)\mathbf 1$ can fail to be regular only at $z=1$ and
$z=-1$.  At $z=-1$ the $\bar M$-conjugation of Lemma~\ref{lem:conj}
turns the problem into the ordinary renewal pole at $+1$ applied to the
test function $(-1)^{E}$, and the condition is
\[
\text{(C2)}\qquad
\bigl\|\kappa_1^{\,k}(-1)^{E}\bigr\|_\infty\le C\rho^k
\quad\text{for some }\rho<1 ,
\]
i.e.\ \emph{the parity of the environment block count decorrelates
geometrically along the post-flip chain}.  At $z=1$ the condition is
$-1\notin\spec\kappa_1$ with a uniform inverse (C1), and we prove that
a single Doeblin minorisation shared by the $n$-th and $(n{+}1)$-st
post-flip laws implies it, with
$\|(I+\kappa_1)^{-1}\|\le(2n+1)/(2\delta)$
(Proposition~\ref{prop:c1}).  Neither condition contains a $z$, a
phase, or an alternating sign that is not the sign of a state function.
(C2) is exactly what \texttt{s11\_kflip} measures, and it holds
numerically with $\rho\approx0.3$ from both a macroscopic and a deeply
trapped start (\S\ref{sec:s11num}).
\end{enumerate}

\subsubsection{The label is a passenger}\label{sec:s11quotient}

Recall the state: a partition of the circle into curves with masses,
two marks $u,v$, and the pair $(p_1,p_2)$ --- the two arcs $(x,y)$ of
the marked curve when $u\sim v$, the two marked curve masses $(a,b)$
when $u\not\sim v$ --- with $s=p_1+p_2$, and $\varepsilon=\pm1$ the
label.  Let $\bar\Gamma=(p_1,p_2;\mathrm{env})$ be the state with the
label forgotten and $\pi_\ast:\Gamma\mapsto\bar\Gamma$ the projection.

\begin{lemma}[label autonomy]\label{lem:autonomy}
The transition kernel of $\Gamma$ projects to a kernel on
$\bar\Gamma$, and in both sectors the moves and their probabilities are
\[
\text{flip: }2p_1p_2,\qquad
p_i\to p_i\sqrt U:\ p_i^2,\qquad
p_i\to p_i+q_j:\ 2p_iq_j,\qquad
\text{environment moves: unchanged,}
\]
a flip replacing $(p_1,p_2)$ by $(V,s-V)$ with
$V\sim\mathrm{U}(0,p_1)\star\mathrm{U}(0,p_2)$ and toggling
$\varepsilon$.  Consequently:
\begin{enumerate}[leftmargin=2em]
\item[(i)] the involution $J$ of the state space that swaps the two
sectors keeping $\bar\Gamma$ fixed commutes with the transition kernel
and fixes every flip time;
\item[(ii)] in the decomposition
$F=g\circ\pi_\ast+\varepsilon\cdot(f\circ\pi_\ast)$ of bounded
functions, $K_z=\kappa_z\oplus(-\kappa_z)$, where
$(\kappa_zg)(\bar x)=\E_{\bar x}[z^\tau g(\bar\Gamma_\tau)]$; hence
$\spec K_z=\spec\kappa_z\cup(-\spec\kappa_z)$ and
$\widehat\Phi(z)=(1-z)^{-1}(I+\kappa_z)^{-1}(I-\kappa_z)\mathbf 1$;
\item[(iii)] if $\bar\pi$ is a stationary law of $\bar\Gamma$ then the
stationary label is uniform: $\Pr[u\sim v]=\tfrac12$.
\end{enumerate}
\end{lemma}

\begin{proof}
A step picks two uniform points of the circle; if they lie in the same
curve it is a split, otherwise a merge.  In the label-$+$ sector the
marked curve has mass $s$ and is picked twice with probability $s^2$;
given that, the two cuts lie on different arcs with probability
$2xy/s^2$ (a flip, separating $u$ from $v$) and on the same arc $i$
with probability $p_i^2/s^2$, in which case the arc loses a factor
$\sqrt U$ in law ($1-|U_1-U_2|\sim\sqrt U$) and the removed piece
becomes an unmarked curve.  A merge of the marked curve with an
environment curve $q_j$ has probability $2sq_j$, and the mass $q_j$ is
attached to arc $i$ with probability $p_i/s$, i.e.\ rate $2p_iq_j$.  In
the label-$-$ sector the same four rates hold with $(a,b)$ in place of
$(x,y)$: picking $u$'s and $v$'s curves is a merge --- a flip, joining
them --- with probability $2ab$; picking $u$'s curve twice splits it,
$a\to a\sqrt U$; picking $u$'s curve and $q_j$ merges them,
$a\to a+q_j$.  The post-flip law is the same trapezoid
$\mathrm{U}(0,p_1)\star\mathrm{U}(0,p_2)$ in both directions
(Lemma~\ref{lem:recharge}(ii)).  This is the displayed list, and it
does not mention $\varepsilon$.  (i) is immediate; (ii) follows because
$\varepsilon(\Gamma_\tau)=-\varepsilon(x)$ while the law of
$\bar\Gamma_\tau$ depends only on $\bar x$; (iii) because $J$ preserves
the stationary law --- it commutes with the kernel and fixes
$\bar\Gamma$ --- while exchanging the two sectors.
\end{proof}

\emph{Status: PROVED.  Implicit in the phrase ``quotient chain'' used
since session 8 but never stated; it is worth stating because it is
what makes the next two subsections possible.  Two consequences.  (a)
The first-flip time $\tau$ has the same law from $x$ and from $Jx$:
the alternating renewal governing $\Phi_d$ is \emph{symmetric}, which
is why $\Phi_d\to0$ rather than to the generic limit
$(\E_+\tau-\E_-\tau)/(\E_+\tau+\E_-\tau)$ of an asymmetric alternating
renewal.  This is the structural reason behind
$\Pr_\pi[\mathrm{tog}]=\tfrac12$, and it is exactly the residue that
has to vanish at $z=1$ below.  (b) All the parity content therefore
lives in $\kappa_z$, a kernel on a space where $\varepsilon$ does not
exist.}

\subsubsection{(R2) is false, in both bookkeepings}\label{sec:s11false}

Write $E(\Gamma)$ for the number of environment curves and let $\bar M$
be multiplication by $(-1)^{E}$, an isometric involution of every
weighted sup-norm space; $N$ denotes multiplication by $\varepsilon$.

\begin{lemma}[sign conjugations]\label{lem:conj}
For every $z\in\mathbb C$,
\[
\text{(i)}\quad N\,K_z\,N=-K_z,
\qquad\qquad
\text{(ii)}\quad \bar M\,\kappa_z\,\bar M=-\kappa_{-z} .
\]
Hence $\spec K_z=-\spec K_z$ and $\spec\kappa_{-z}=-\spec\kappa_z$, and
the resolvents $(I\pm\kappa_{\pm z})^{-1}$ have equal norms in pairs.
\end{lemma}

\begin{proof}
(i) Exactly one flip has occurred at time $\tau$, so
$\varepsilon(\Gamma_\tau)=-\varepsilon(x)$ a.s.
(ii) By Lemma~\ref{lem:parity}, $(-1)^{N_t}=(-1)^t(-1)^{E_t-E_0}$; at
$t=\tau$ we have $N_\tau=1$, so $(-1)^\tau=-(-1)^{E_\tau-E_0}$.
Therefore
$(\bar M\kappa_z\bar Mg)(\bar x)=(-1)^{E(\bar x)}
\E_{\bar x}\bigl[z^\tau(-1)^{E(\bar\Gamma_\tau)}g(\bar\Gamma_\tau)\bigr]
=-\E_{\bar x}\bigl[(-z)^\tau g(\bar\Gamma_\tau)\bigr]$.
\end{proof}

\emph{Remark (the label-carrying form of (ii)).  Since $B$ changes by
exactly $\pm1$ at every step, $(-1)^\tau=(-1)^{B(\Gamma_\tau)-B(x)}$
directly, so with $M$ = multiplication by $(-1)^{B}$ one has the
cleaner-looking $M K_zM=K_{-z}$ on the label-carrying space, and
$K_{-1}[(-1)^B]=+(-1)^B$.  That eigenvalue is $+1$ and is harmless;
the obstruction to (R2) is the eigenvalue $-1$ of
$\kappa_{-1}$ at $(-1)^E$ in Proposition~\ref{prop:r2false}.  The two
are consistent: $(-1)^B=-\varepsilon\,(-1)^E$, and
$N K_zN=-K_z$ supplies the sign.  Keeping these apart matters: in the
label-carrying picture $\E_x[(-1)^{\tau_k}]$ has no a-priori reason to
decay, while in the label-free picture its decay is exactly (C2).}

\begin{proposition}[failure of (R2)]\label{prop:r2false}
Assume $\tau<\infty$ a.s.\ (as is taken for granted after
Proposition~\ref{prop:genfun}), so that $\kappa_1$ is Markov.  Then
\[
K_1\varepsilon=-\varepsilon
\qquad\text{and}\qquad
\kappa_{-1}\bigl[(-1)^{E}\bigr]=-(-1)^{E} .
\]
Hence $I+K_z$ is not injective at $z=1$ in the label-carrying
bookkeeping and $I+\kappa_z$ is not injective at $z=-1$ in the
label-free one, on every function space containing the bounded
functions.  In particular condition (R2) of
Proposition~\ref{prop:condecay} never holds, and
Proposition~\ref{prop:condecay} is vacuous as stated.
\end{proposition}

\begin{proof}
The first identity is Lemma~\ref{lem:conj}(i) applied to $\mathbf 1$
together with $K_1\mathbf 1=\mathbf 1$.  For the second, by
Lemma~\ref{lem:conj}(ii) with $z=1$,
$\kappa_{-1}\bar M\mathbf 1=-\bar M\kappa_1\mathbf 1=-\bar M\mathbf 1$.
\end{proof}

\emph{Status: PROVED.  The failure is structural, not technical: the
null vectors $\varepsilon$ and $(-1)^E$ are bounded by $1$ and lie in
every space in sight, so no choice of weight, state space or test class
removes them.  Session 10 did not see it because
Proposition~\ref{prop:genfun} applies $(I+K_z)^{-1}$ only to
$(I-K_z)\mathbf 1$, which degenerates at the same points; the whole
question is whether the vanishing of the numerator beats the blow-up of
the resolvent, and that is not what (R2) says.  The two failures are
the same failure in the two bookkeepings, since
$\spec K_z=\pm\spec\kappa_z$ by Lemma~\ref{lem:autonomy}(ii).}

\begin{proposition}[the two singular points, exactly]\label{prop:repair}
For $|z|<1$,
$\widehat\Phi(z)=(1-z)^{-1}(I+\kappa_z)^{-1}(I-\kappa_z)\mathbf 1$, and
for every $w$,
\[
(I+\kappa_{-w})^{-1}(I-\kappa_{-w})\mathbf 1
=\bar M\,(I-\kappa_{w})^{-1}(I+\kappa_{w})\,\bar M\mathbf 1 .
\]
Thus near $z=1$ the singular factor is $(I+\kappa_z)^{-1}$ acting on
the numerator $(I-\kappa_z)\mathbf 1=\E_\cdot[1-z^\tau]=O(|1-z|)$,
while near $z=-1$ it is the ordinary renewal resolvent
$(I-\kappa_w)^{-1}$, $w\to1$, acting on
$(I+\kappa_w)(-1)^{E}\to2\,(-1)^{E}$.
\end{proposition}

\begin{proof}
The first statement is Lemma~\ref{lem:autonomy}(ii) substituted in
Proposition~\ref{prop:genfun}; the second is Lemma~\ref{lem:conj}(ii)
together with $\bar M^2=I$.
\end{proof}

\emph{Status: PROVED (algebra).  This is the sharp form of the
correction.  At $z=-1$ the singularity is the classical simple pole of
a Markov renewal resolvent at $w=1$, whose residue is the stationary
projection $\langle\bar\pi,\cdot\rangle\mathbf 1$; the residue of
$\widehat\Phi$ at $z=-1$ is therefore proportional to
$\langle\bar\pi,(-1)^{E}\rangle$, and $\Phi_d$ carries a non-decaying
alternating part $c\,(-1)^d$ unless that pairing vanishes.  That is
precisely the alternation found in Lemma~\ref{lem:parity}(a) --- where
it is measured to decay like $(-1)^dc_1/d^2$, i.e.\ the residue does
vanish.  At $z=1$ the numerator vanishes and the residue to be beaten
is the residue of a \emph{symmetric} alternating renewal, zero by
Lemma~\ref{lem:autonomy}(a).  Both residues therefore vanish for
structural reasons; what is unproved is the quantitative, uniform
version, and that is (C1) and (C2) of \S\ref{sec:s11left}.}

\subsubsection{The idle loop: an unconditional contraction off
\texorpdfstring{$z=\pm1$}{z=+-1}}\label{sec:idleloop}

Write the \emph{arc-refined mass list} $r(\Gamma)=(r_j)_j$ for the
masses of the unmarked curves together with $p_1$ and $p_2$; so
$\sum_jr_j=1$ and $r$ has at most $B(\Gamma)+1$ entries, $B$ the total
number of curves.

\begin{lemma}[the idle loop]\label{lem:idleloop}
Let $L$ be the event that step~1 splits some curve and step~2 merges
the two pieces that step~1 created, the returned piece being
re-attached to the arc it came from if the split curve was the marked
one in the label-$+$ sector.  On $L$ no flip occurs and
$\bar\Gamma_2=\bar\Gamma_0$; up to a null set
$L=\{\bar\Gamma_2=\bar\Gamma_0\}$; and
\[
\lambda(\Gamma):=\Pr_\Gamma[L]\;=\;\frac13\sum_j r_j(\Gamma)^4
\;\ge\;\frac1{3\,(B(\Gamma)+1)^3}.
\]
\end{lemma}

\begin{proof}
Fix an entry $r$ of the arc-refined list; the events for distinct
entries are disjoint.

\emph{$r=q$, an unmarked curve.}  Step~1 picks it twice with
probability $q^2$ and cuts it at two independent uniform points,
giving pieces $qV$ and $q(1-V)$ with $V=|U_1-U_2|$, density $2(1-v)$.
Step~2 picks those two curves with probability $2\,qV\cdot q(1-V)$ and
merges them back into a curve of mass $q$.  Total
$q^2\cdot2q^2\,\E[V(1-V)]=q^4/3$, since
$\E[V(1-V)]=\int_0^1v(1-v)\,2(1-v)\,dv=\tfrac16$.

\emph{$r=a=p_1$ in the label-$-$ sector.}  Identical, the mark staying
in the piece of mass $a(1-V)$: total $a^4/3$.

\emph{$r=x$, one arc in the label-$+$ sector.}  Both cuts must fall in
that arc, probability $x^2$ --- and this is a split, not a flip,
precisely because a flip needs the cuts on \emph{different} arcs.  The
piece has mass $w=xV$; the marked curve keeps $s-w$ and the arc becomes
$x-w$.  Step~2 merges the piece back with probability $2w(s-w)$, and
the returned mass is re-attached to the arc it came from with
conditional probability $(x-w)/(s-w)$.  The factor $s-w$ cancels:
$x^2\,\E[2\,xV\cdot x(1-V)]=2x^4\E[V(1-V)]=x^4/3$.

Summing gives $\lambda=\tfrac13\sum_jr_j^4$.  Conversely a merge at
step~1 followed by a split at step~2 reproduces the two merged masses
with probability $0$, and a split followed by the merge of two other
curves reproduces the mass list only on a null set.  Finally, with
$p\le B+1$ entries, Cauchy--Schwarz twice gives
$\sum r_j^4\ge(\sum r_j^2)^2/p\ge p^{-3}$.
\end{proof}

\emph{Status: PROVED.  VERIFIED (\texttt{s11\_ztau}, $2\times10^6$
two-step trials per start, exact state comparison at tolerance
$10^{-12}$): one curve, arcs $(\tfrac12,\tfrac12)$: $\lambda$ measured
$0.041530$ vs.\ $\tfrac13\cdot2\cdot2^{-4}=0.041667$; $s_0=0.6$, arcs
$(0.1,0.5)$, environment $0.4$: $0.029459$ vs.\ $0.029400$;
$s_0=0.5$, arcs $(0.02,0.48)$, environment $0.5$: $0.038556$ vs.\
$0.038528$.  The formula is an identity, not an estimate, and it is the
\emph{arc-refined} list that appears: splitting the marked curve across
its two arcs is a flip, not an idle loop.  The quantity $\sum_jr_j^4$
is the one whose expectation Theorem~\ref{thm:sizebiased} computes
exactly, $\E\sum p^4(m)=\tfrac14+\int_0^1(1-x^3)(1-2x)^mdx\to\tfrac14$,
so $\E\lambda\to\tfrac1{12}$; the block bound is needed only to turn
that average into a uniform bound.}

\begin{theorem}[idle-loop contraction]\label{thm:idleloop}
Let $\lambda_*=\inf_{\bar x}\lambda(\bar x)$ over the state space on
which $\kappa_z$ acts.  For all $|z|\le1$,
\[
\|\kappa_z\|_{\infty\to\infty}\;\le\;
\sup_{\bar x}\frac{1-\lambda(\bar x)}{|1-\lambda(\bar x)z^2|}
\;\le\;\frac1{1+\lambda_*\,(1-\Re z^2)} .
\]
Consequently if $1-\Re z^2\ge t_0$ then $I\pm\kappa_z$ and
$I-\kappa_z^2$ are invertible with
$\|(I\pm\kappa_z)^{-1}\|\le1+(\lambda_*t_0)^{-1}$.  On the state space
$\Omega_K=\{B\le K\}$ --- the chain killed at the first split that
would make $B=K+1$, whose $\Phi_d$ differs from the true one by at most
$\Pr[H^c]\le(9d+1)^{-\gamma}$ (Corollary~\ref{cor:horizon}) --- one has
uniformly
\[
\|\kappa_z\|_\infty\;\le\;1-\min\Bigl(\frac{t_0}{6K^3},\ \frac1{2K}\Bigr),
\qquad\text{so}\qquad
\|(I\pm\kappa_z)^{-1}\|\le\max\Bigl(\frac{6K^3}{t_0},\,2K\Bigr).
\]
\end{theorem}

\begin{proof}
$L$ is measurable with respect to the first two steps; on $L$ no flip
occurs and $\bar\Gamma_2=\bar x$.  By the Markov property at time $2$,
$\E_{\bar x}[z^\tau g(\bar\Gamma_\tau)\mathbf 1_L]
=z^2\lambda(\bar x)(\kappa_zg)(\bar x)$, whence
\[
(1-\lambda z^2)(\kappa_zg)(\bar x)
=\E_{\bar x}\bigl[z^\tau g(\bar\Gamma_\tau)\mathbf 1_{L^c}\bigr],
\qquad
|(\kappa_zg)(\bar x)|\le\frac{(1-\lambda)\|g\|_\infty}{|1-\lambda z^2|} .
\]
For $|z|\le1$, $|1-\lambda z^2|\ge\Re(1-\lambda z^2)
=(1-\lambda)+\lambda t$ with $t=1-\Re z^2$, and
$u\mapsto(1-u)/(1-u+ut)$ has derivative $-t/(\,\cdot\,)^2<0$, so the
supremum over $\bar x$ is attained at $\lambda_*$ and equals
$1/(1+\lambda_*t/(1-\lambda_*))\le1/(1+\lambda_*t)$.  The resolvent
bounds are Neumann series.

For the killed chain, split the states.  If $B(\bar x)\le K-1$ the idle
loop raises $B$ only to $B+1\le K$, so it stays inside $\Omega_K$ and
is not killed; Lemma~\ref{lem:idleloop} gives
$\lambda\ge\tfrac13K^{-3}$, and $1/(1+u)\le1-u/2$ for $u\le1$ gives the
first branch.  If $B(\bar x)=K$ then \emph{every} split kills, and a
split occurs at the next step with probability
$\sum_i(\text{curve mass})^2\ge1/K$, so
$|\kappa_zg(\bar x)|\le\|g\|_\infty\Pr_{\bar x}[\tau<\text{kill}]
\le(1-1/K)\|g\|_\infty$.  Take the worse of the two branches.
\end{proof}

\emph{Status: PROVED.  VERIFIED (\texttt{s11\_ztau}): from each of the
three deterministic starts above,
$|\E_{\bar x}[z^\tau]|=|\kappa_z\mathbf 1(\bar x)|$ was measured at
$\theta/\pi=0.1,\dots,0.9$ ($2\times10^5$ trials each) and compared
with $(1-\lambda)/|1-\lambda z^2|$: the bound holds at every point,
with margin (at $\theta=\pi/2$ from the $(\tfrac12,\tfrac12)$ start,
measured $0.446$ against the bound $0.920$).  The bound is crude --- it
uses one two-step return and nothing else --- but it is \emph{uniform
in the state}, which is the entire difficulty; and it is the right
mechanism, since $1-\lambda z^2$ is exactly the Euler factor of the
geometric component that the idle loop puts into the law of $\tau$ at
lag $2$, which is what an aperiodicity or non-lattice hypothesis is
normally assumed to supply.}

\emph{Three items on the open list are settled by
Theorem~\ref{thm:idleloop} and Lemma~\ref{lem:conj}.  (i) Pitfall~45
and TODO~11b(v) --- the non-lattice condition ``on the whole circle,
not only near $z=1$'' --- are PROVED, with rate $1-\Re z^2$ and
constant $3K^3$.  (ii) TODO~11b(iv), ``$-1\notin\spec K_1$'', is
DISPROVED and was the wrong question
(Proposition~\ref{prop:r2false}).  (iii) The point $z=-1$, treated as a
separate problem in every handoff since session 8, is by
Proposition~\ref{prop:repair} the ordinary renewal point $w=1$ with the
test function $(-1)^E$: not a separate problem, and not a place where
anything new has to be invented.}

\subsubsection{What is left: (C1) and (C2)}\label{sec:s11left}

By Theorem~\ref{thm:idleloop}, $\widehat\Phi$ is regular on
$\{|z|\le1\}\setminus\{\pm1\}$ with bounds polynomial in $K$.  Two
conditions remain, both about the single Markov kernel $\kappa_1$ ---
the post-flip quotient chain --- on $\Omega_K$:

\medskip
\noindent\textbf{(C1)}\quad$-1\notin\spec\kappa_1$, with
$\|(I+\kappa_1)^{-1}\|\le C(K)$;\\
\noindent\textbf{(C2)}\quad$\bigl\|\kappa_1^{\,k}(-1)^{E}\bigr\|_\infty
\le C(K)\rho^k$ for some $\rho=\rho(K)<1$.
\medskip

\noindent With (R1) and these two, the proof of
Proposition~\ref{prop:condecay} goes through: (C1) gives regularity at
$z=1$, where the numerator has its zero; (C2) gives, through
Proposition~\ref{prop:repair}, both the vanishing of the residue at
$z=-1$ --- since $\langle\bar\pi,(-1)^E\rangle=\lim_k\kappa_1^k(-1)^E=0$
--- and a Hölder modulus there.  \emph{Neither contains a $z$, a phase,
or an alternating sign that is not the sign of a state function: the
parity content of the problem has been reduced to the single statement
that the environment block count has no asymptotic parity bias along
the post-flip chain.}

\emph{Neither (C1) nor (C2) is false for the reason (R2) was.  The
only deterministically-toggled state function available on the quotient
is $(-1)^E$ (flips leave $E$ fixed; every non-flip step changes it by
$\pm1$), and by Lemma~\ref{lem:conj}(ii)
$\kappa_1[(-1)^E]=-(-1)^{E}\,\E_\cdot[(-1)^\tau]$, so $(-1)^E$ is an
eigenfunction of $\kappa_1$ if and only if $\E_x[(-1)^\tau]$ is
constant in $x$ --- which it measurably is not
($-0.1306\pm0.003$ from the macroscopic start against
$-0.0009\pm0.005$ from the trapped start, \S\ref{sec:s11num}).}

\begin{proposition}[a minorisation suffices for (C1)]\label{prop:c1}
Suppose there are $n\ge1$, $\delta>0$ and a probability measure $\nu$
with
\[
\Pr_{\bar x}\bigl[\bar\Gamma_{\tau_n}\in\cdot\,\bigr]\ge\delta\nu(\cdot)
\qquad\text{and}\qquad
\Pr_{\bar x}\bigl[\bar\Gamma_{\tau_{n+1}}\in\cdot\,\bigr]\ge\delta\nu(\cdot)
\qquad\text{for every }\bar x .
\]
Then $I+\kappa_1$ is invertible on $L^\infty$ with
$\|(I+\kappa_1)^{-1}\|\le(2n+1)/(2\delta)$.
\end{proposition}

\begin{proof}
Let $(I+\kappa_1)F=G$ and $m=\|F\|_\infty$.  Iterating
$F=G-\kappa_1F$ gives $F=H_j+(-1)^j\kappa_1^{\,j}F$ with
$H_j=\sum_{i<j}(-1)^i\kappa_1^{\,i}G$, $\|H_j\|\le j\|G\|$.  Write
$\kappa_1^{\,n}F=\delta c+(1-\delta)\rho^1_{\bar x}(F)$ and
$\kappa_1^{\,n+1}F=\delta c+(1-\delta)\rho^2_{\bar x}(F)$ with the same
$c=\nu(F)$ and probability measures $\rho^i_{\bar x}$.  Multiplying the
$j=n$ identity by $(-1)^n$ and the $j=n+1$ identity by $(-1)^{n+1}$ and
subtracting the results eliminates the $\delta c$ terms' difference and
leaves
\[
2F(\bar x)=\pm\bigl(H_n+H_{n+1}\bigr)(\bar x)
+(1-\delta)\bigl(\rho^2_{\bar x}-\rho^1_{\bar x}\bigr)(F),
\]
so $2m\le(2n+1)\|G\|+2(1-\delta)m$, i.e.\
$2\delta m\le(2n+1)\|G\|$.
\end{proof}

\emph{The proof uses no property of $\kappa_1$ beyond
$\|\kappa_1\|\le1$: it is valid verbatim for the \emph{sub}-Markov
$\kappa_1$ of the chain killed on $\Omega_K$, the residual measures
$\rho^i_{\bar x}$ then having mass $\le1-\delta$ rather than exactly
$1-\delta$.  This matters because $\Omega_K$ is where
Theorem~\ref{thm:idleloop} lives.}

\emph{\textbf{SESSION-12 CORRECTION.}  The hypothesis of this proposition
is never satisfied: for every $n$ there is no common minorant of the
$n$-th post-flip laws over $\Omega_K$, because a curve of mass
$\varepsilon$ survives the first $n$ flips untouched with probability
$\to1$ (Proposition~\ref{prop:nomino}).  The proposition stays correct
and vacuous; TODO~12b is closed.}

\emph{Status: PROVED as a reduction; the minorisation is OPEN.  It is
the honest remaining difficulty and it is the classical one: the
post-flip chain regenerates on the fibre --- two consecutive flips
produce a $p_1''$ whose law dominates
$\tfrac14\mathrm{Unif}[s/4,3s/4]$ (\S\ref{sec:markovrenewal}) --- but
not in the base $(s,\mathrm{env})$, which is carried along.  On
$\Omega_K$ the base has at most $K$ curves, and the numerics of
\S\ref{sec:s11num} show it equilibrating in about four flips from a
deliberately bad start, so a minorisation with $n=O(1)$ and a
$\delta$ polynomial in $1/K$ is the natural target; the $K$-dependence
of $\delta$ may be as bad as one likes, polynomially.}

\emph{Routes to (C2).  (a) Uniform ergodicity plus a parity symmetry:
the same minorisation gives
$\|\kappa_1^k(F-\bar\pi F)\|\le C\rho^k$ for all bounded $F$, reducing
(C2) to $\langle\bar\pi,(-1)^E\rangle=0$ --- the statement that the
post-flip stationary law on $\Omega_K$ has no environment block-count
parity bias.  On $\Omega_K$ that is a statement about a
finite-dimensional chain and is the cleanest target for session 12.
(b) Directly: by Lemma~\ref{lem:conj}(ii),
$\kappa_1^{\,k}(-1)^E=(-1)^{E}(-1)^k\kappa_{-1}^{\,k}\mathbf 1$ and
$\kappa_{-1}^{\,k}\mathbf 1(\bar x)=\E_{\bar x}[(-1)^{\tau_k}]$, so
(C2) says exactly that the parity of the $k$-th flip time decorrelates
geometrically --- which is what \texttt{s11\_kflip} measures.  A proof
along (b) needs a state-uniform lower bound on the probability that
two consecutive inter-flip times have opposite parity; the idle loop
supplies a uniform event at lag $2$, but it changes $\tau$ by an
\emph{even} amount, so a genuinely odd perturbation is still needed.
That is the one place where an odd/even asymmetry must be produced
rather than conjugated away, and it is the sharpest statement of the
residue of the whole problem.}

\emph{A Lyapunov remark.}  The proof of Theorem~\ref{thm:idleloop}
degrades only through $\lambda^{-1}=3(\sum_jr_j^4)^{-1}$, whose
expectation along the chain is known exactly
(Theorem~\ref{thm:sizebiased}).  A drift inequality
$\kappa_1V\le\varrho V+b$ with $V=(\sum_jr_j^4)^{-\theta}$, $\theta>0$
small, would replace the killing at $B=K$ by a weighted space and
supply the ``base'' half of the minorisation of
Proposition~\ref{prop:c1} at the same time.  It is a finite
computation of the same kind as Lemma~\ref{lem:logdrift} and is the
recommended first task of session 12.

\subsubsection{Numerics along the post-flip chain}\label{sec:s11num}

TODO~11d asked for $\E_{\bar x}[(-1)^{\tau_k}]$ along the post-flip
chain, predicting geometric decay under (R2).  The prediction needed
correcting before the run could be interpreted --- under (R2) as stated
there is no such prediction, (R2) being false --- but the quantity is,
by route (b) above, \emph{exactly} the residual condition (C2).
Measured (\texttt{s11\_kflip}: stationary environment burnt in over
$4096$ steps, accepted on the initial gap, $k\le16$, no censoring):

\emph{Macroscopic start, gap in $[\tfrac18,\tfrac38)$,
$1.2\times10^5$ samples (s.e.\ $0.0029$):
$\E[(-1)^{\tau_k}]=-0.1306,\ +0.0263,\ -0.0100,\ +0.0011$, and at the
noise floor for every $k\ge4$: geometric decay with ratio
$\approx0.2$--$0.4$ and no persistent $(-1)^k$ component.  Trap start,
gap in $[10^{-3},10^{-2})$, $4\times10^4$ samples (s.e.\ $0.0050$):
$-0.0009,\,+0.0018,\,+0.0024,\dots$, at the noise floor throughout ---
one long inter-flip time randomises the parity immediately.
\textbf{This is (C2), measured, with $\rho\approx0.3$.}}

\emph{\textbf{SESSION-12 CORRECTION.}  The ``$\rho\approx0.3$'' is the
transient.  With $2\times10^7$ samples from fixed low-block-count
starts, $a_k=(-1)^k(2.5\pm0.1)k^{-3}$ for $7\le k\le12$: the decay is
polynomial, not geometric, and the stationary-start average above
hides it because it averages the sign $(-1)^{B(\bar x)}$ over the
start (\S\ref{sec:s12num}).  (C2) is replaced by the polynomial
(C2$''$) of \S\ref{sec:s12left}, which suffices for a H\"older exponent
$(p-1)/(p+1)>4/9$ at $z=-1$ when $p>13/5$.}

\emph{Inter-flip parity $\E[(-1)^{\tau_k-\tau_{k-1}}]$ settles at
$-0.170\pm0.003$ for $k\ge3$ from both starts: bounded away from
$\pm1$, as (C2) needs, and strikingly stable.}

\emph{The post-flip chain equilibrates in about four flips.  From the
trap start, $\E[\mathrm{gap}]$ at $\tau_k$ runs
$0.1995,0.2375,0.2468,0.2484,0.2491,\dots$ and $\E[s]$ runs
$0.715,0.774,0.790,0.797,0.798$, matching the macroscopic start's
$0.2498$ and $0.799$ by $k=4$; the correlations
$\E[\mathrm{sgn}(g_0-\tfrac14)\,\mathrm{sgn}(g_k-\tfrac14)]$ decay
$0.317,0.090,0.022,0.012,\dots$, geometrically with ratio
$\approx0.3$.  This is direct evidence for the minorisation of
Proposition~\ref{prop:c1} with $n$ of order $4$.}

\emph{First-flip tails along the chain: $t^2\Pr[\tau_k-\tau_{k-1}>t]$
lies in $15$--$20$ on $16\le t\le128$ for $k=2,4,8$ from both starts
(and $\approx10^4$ for $k=1$ from the trap start, consistent with the
$g^{-2}$ prefactor), so (R1) holds uniformly \emph{along} the post-flip
chain and not merely at the start: TODO~11d(iii), answered in the
affirmative.}

\emph{Finally, \texttt{s11\_kflip} asserts
$B(\Gamma_t)-B(\Gamma_0)\equiv t\pmod2$ at every flip of every one of
the $1.6\times10^5$ trajectories --- the hypothesis of
Lemma~\ref{lem:conj}(ii), via Lemma~\ref{lem:parity} --- and it holds.}

%----------------------------------------------------------------------
\subsection{Session 12: the minorisation route is closed by inert dust;
(C2) is polynomial, not geometric}\label{sec:s12}

Session 12 was directed at (C1)+(C2) of \S\ref{sec:s11left}, through
TODO~12a (a Lyapunov drift for $\lambda$), 12b (the minorisation of
Proposition~\ref{prop:c1}) and 12c (an odd perturbation), with the
instruction to return one of (A1)--(A3) or a clean negative (B).  The
outcome is (B), in a form sharper than the one the brief anticipated.
Four findings, in logical order.

\begin{enumerate}[leftmargin=2.2em]
\item[(1)] \textbf{The environment carries inert dust, and this closes
every Doeblin-type route.}  A curve of mass $\varepsilon$ is touched by
a chord with probability $\le2\varepsilon$ per step, so it survives
\emph{unchanged} through the first $n$ flips with probability
$\to1$ as $\varepsilon\to0$ (Lemma~\ref{lem:inert}).  Consequently the
post-flip laws from the states $x_\varepsilon$ (a fixed macroscopic
marked pair, plus one dust curve of mass $\varepsilon$) are supported,
up to $o(1)$, on the pairwise disjoint sets ``some curve has mass
exactly $\varepsilon$''.  No common minorant exists: the hypothesis of
Proposition~\ref{prop:c1} fails on $\Omega_K$ for \emph{every} $n$
(Proposition~\ref{prop:nomino}), and so does uniform ergodicity of
$\kappa_1$ in $L^\infty(\Omega_K)$ and $V$-uniform ergodicity for
\emph{every} weight $V$ (Corollary~\ref{cor:noerg}).  Route (a) of
\S\ref{sec:s11left} and TODO~12b are therefore closed, not
deferred.  The proof uses only $\tau_n<\infty$ a.s.
\item[(2)] \textbf{TODO~12a fails on the full state space.}  From a
state whose non-dust material has total mass $m$ and whose remaining
mass sits in $N\gg m^{-3}$ equal pieces, one has
$\sum_jr_j^4\le17m^4$ deterministically for a long time, so
$\kappa_1^kV\ge\tfrac34(17m^4)^{-\theta}$ while $V\le m^{-4\theta}$:
no inequality $\kappa_1V\le\varrho V+b$ with $\varrho<1$ can hold
(Proposition~\ref{prop:nodrift}).  As instructed, the fallback is the
killed chain on $\Omega_K$, where $V$ is bounded anyway.
\item[(3)] \textbf{(C2) is true but polynomial.}  With $2\times10^7$
samples from the one-curve start,
$a_k(x)=\E_x[(-1)^{\tau_k}]=(-1)^k\,(2.5\pm0.1)\,k^{-3}$ for
$7\le k\le13$, and with $10^8$ samples the local exponent over
$8\le k\le20$ is $2.7\pm0.3$; a tail of the same form,
$k^3|a_k|\to2.0$--$2.4$, from the three-curve start $(0.3,0.3;\{0.4\})$,
and the same tail \emph{emerging} after a transient of $\approx8$
flips from a nine-curve start.  Session~11's ``$\rho\approx0.3$''
was the transient above a noise floor of $0.003$.  Since
$a_k=\pm\E_x[(-1)^{B(\bar\Gamma_{\tau_k})}]$ is the parity bias of the
block count at the $k$-th flip, the mechanism is the slow ($\log$)
growth of $B$: the block-count parity is a slow mode of $\kappa_1$ at
eigenvalue $+1$, not a mode near $-1$ (\S\ref{sec:s12num}).  The
polynomial decay is insensitive to inserting dust into the start
(identical to two standard errors) and is \emph{weaker} from starts
with $9$ or $17$ curves.
\item[(4)] \textbf{What the reduction actually needs is polynomial.}
At $z=-1$ Proposition~\ref{prop:condecay} needs only
$G\in C^{0,\gamma}$, $\gamma>4/9$, and a polynomial decay of the
twisted parity $\E_x[w^{\tau_k}(-1)^{E_{\tau_k}}]$ with exponent $p$,
uniformly for $w$ near $1$, gives $\gamma=(p-1)/(p+1)$ by a crude
bound (Proposition~\ref{prop:holder}), so the budget needs $p>13/5$;
if the parity memory factorises from the time weighting, as it does
numerically, the exponent improves to $\min(1,p-1)$ and the budget
needs only $p>13/9$.  The measured $p=2.7\pm0.3$ meets the second
comfortably and the first marginally.  (C2) is accordingly replaced by
(C2$''$) in \S\ref{sec:s12left}.  Any proof must be \emph{dust-blind},
i.e.\ work in a topology in which $x$ and $x_\varepsilon$ are close;
we record why the natural transport metric is not contracted by the
dynamics and why an adversarial environment cannot be allowed
(\S\ref{sec:s12left}).
\end{enumerate}

\subsubsection{Inert dust}\label{sec:s12dust}

Throughout, the chain is realised on the circle: the state is a
partition of the circle into curves (each a cyclically ordered union of
arcs, the cyclic order being that of the chord diagram), a step picks
two independent uniform points, and the move is determined by the
curves the points lie in and their positions along those curves.

\begin{lemma}[inert dust coupling]\label{lem:inert}
Let $y$ be a state, $D$ a sub-arc of the circle of mass $\varepsilon$
lying inside one curve of $y$ and containing neither mark, and let
$y_\varepsilon$ be the state obtained from $y$ by declaring $D$ a
separate (unmarked) curve.  Drive the chains from $y$ and from
$y_\varepsilon$ by the same uniform points.  Then, up to the first step
at which a point falls in $D$, the state of the $y_\varepsilon$-chain
is the state of the $y$-chain with $D$ declared a separate curve; in
particular the two chains have the same flip times, and the
$y_\varepsilon$-chain contains a curve of mass exactly $\varepsilon$.
The first point in $D$ falls at a step of probability at most
$2\varepsilon$, so for every $T$,
\[
\Pr\bigl[\text{the coupling holds through time }T\bigr]\ge1-2\varepsilon T .
\]
\end{lemma}

\begin{proof}
$D$ is a contiguous segment of the cyclic order of the curve $C\supset D$
containing it.  A split of $C$ at two points outside $D$ cuts the cycle
into two arcs, one of which contains $D$; in the $y_\varepsilon$-chain
the same cuts split $C\setminus D$ into the same two arcs with $D$
excised.  A merge of $C$ with another curve at points outside $D$
splices the two cycles at those points; in the $y_\varepsilon$-chain the
same splice is applied to $C\setminus D$.  Moves not involving $C$ are
identical.  A flip is a split or merge decided by which marked material
the two points lie in; the marked material of the two chains differs
only by $D$, which contains no point.  Induction on the steps.
\end{proof}

\begin{proposition}[no state-uniform minorisation]\label{prop:nomino}
Let $K\ge4$ and assume $\tau_n<\infty$ a.s.\ from the state
$x_0=(\tfrac14,\tfrac14;\{\tfrac12\})$ (apart sector).  Then there are
no $n\ge1$, $\delta>0$ and probability measure $\nu$ on $\Omega_K$ with
\[
\Pr^{\Omega_K}_{\bar x}\bigl[\bar\Gamma_{\tau_n}\in\cdot\,\bigr]\ge\delta\,\nu(\cdot)
\qquad\text{for every }\bar x\in\Omega_K\cap S'_K ,
\]
for the chain killed on $\Omega_K$ or for the unkilled chain.  A
fortiori the hypothesis of Proposition~\ref{prop:c1} is never
satisfied.
\end{proposition}

\begin{proof}
Suppose such $n,\delta,\nu$ exist.  For $i\ge3$ let
$\varepsilon_i=2^{-i}$, let $x_i$ be $x_0$ with a sub-arc $D_i$ of mass
$\varepsilon_i$ of its environment curve declared a separate curve
(so $x_i=(\tfrac14,\tfrac14;\{\tfrac12-\varepsilon_i,\varepsilon_i\})$,
$B(x_i)=4$, $x_i\in\Omega_K\cap S'_K$), and let
$S_i=\{\bar x\in\Omega_K:\ \text{some curve of }\bar x\text{ has mass
exactly }\varepsilon_i\}$.  Apply the minorisation at $x_i$ to the set
$\Omega_K\setminus S_i$.  By Lemma~\ref{lem:inert} with $T=T_i$, on the
event that the coupling holds through $T_i$ and $\tau_n\le T_i$ for the
unkilled $x_0$-chain, the $x_i$-chain performs its $n$-th flip at the
same time and, if it is still alive, its state then contains $D_i$
intact; if it has been killed, its state is the cemetery, which is not
in $\Omega_K\setminus S_i$.  Hence
\[
\delta\,\nu(\Omega_K\setminus S_i)\;\le\;2\varepsilon_iT_i+\Pr_{x_0}[\tau_n>T_i]
\qquad\text{for every }T_i ,
\]
and $T_i=\varepsilon_i^{-1/2}$ gives $\nu(S_i)\to1$ as $i\to\infty$.
But a configuration in $\Omega_K$ has at most $K$ curves, hence at most
$K$ distinct masses, hence lies in at most $K$ of the sets $S_i$; so
$\sum_i\nu(S_i)=\int\#\{i:\bar x\in S_i\}\,\nu(d\bar x)\le K<\infty$,
contradicting $\nu(S_i)\to1$.
\end{proof}

\emph{Status: PROVED.  Note what the proof does not use: no moment of
$\tau$, no property of the environment beyond the trivial one that a
curve of mass $\varepsilon$ is rarely hit, and no bound on $K$ beyond
$K\ge4$.  It also shows that the difficulty is not the one anticipated
in TODO~12b (``$K^{-2K}\to K^{-O(1)}$''): the obstruction is not the
\emph{cost} of driving the environment to a reference configuration but
the fact that no finite sequence of steps can remove a curve that is
never hit.  The same applies with the $(n,n+1)$ pair of
Proposition~\ref{prop:c1}, to any weighted space, and to any
``small set'' containing a state with a dust curve --- which every set
of positive $\bar\pi$-measure does, since dust is created at every
scale at rate $\asymp2\eta\,d\eta$ per step.}

\begin{corollary}[no uniform or $V$-uniform ergodicity]\label{cor:noerg}
Let $K\ge4$ and let $\kappa_1$ be the post-flip kernel of the unkilled
chain.
\begin{enumerate}[leftmargin=2em]
\item[(i)] There are no probability measure $\bar\pi$, $C<\infty$ and
$\rho<1$ with $\sup_{\bar x\in\Omega_K}|\kappa_1^kF(\bar x)-\bar\pi F|
\le C\rho^k\|F\|_\infty$ for all bounded $F$ and all $k$.
\item[(ii)] Let $\mathcal Y=\{(a',\tfrac12-a';\{\tfrac12\}):
a'\in[\tfrac18,\tfrac38]\}$ and assume
$(\mathrm M_p)$: $\sup_{y\in\mathcal Y}\E_y[\tau_k]\le Ck^p$ for some
$C,p<\infty$.  Then there is no $V\ge1$ for which $\kappa_1$ is
$V$-uniformly ergodic: no $\bar\pi$, $C'$, $\rho<1$ with
$|\kappa_1^kF(\bar x)-\bar\pi F|\le C'\rho^kV(\bar x)\|F\|_V$ for all
$F$ with $\|F\|_V=\sup|F|/V<\infty$ and all $\bar x$ with $B(\bar x)\le4$.
\end{enumerate}
\end{corollary}

\begin{proof}
(i) With $F_i=\mathbf 1_{S_i}$ (notation of the previous proof) and
$x_i$ as there, Lemma~\ref{lem:inert} gives
$\kappa_1^kF_i(x_i)\ge1-2\varepsilon_iT-\Pr_{x_0}[\tau_k>T]$, while
$\kappa_1^kF_i(x_0)=0$ because from $x_0$ every curve mass at every time
is either a sum of some of $\tfrac14,\tfrac14,\tfrac12$ or has an
atomless law, so no curve ever has mass exactly $\varepsilon_i\le\tfrac18$ a.s.  Let
$i\to\infty$: the oscillation of $\kappa_1^kF_i$ over $\Omega_K$ tends
to $1$ for each fixed $k$, contradicting $2C\rho^k<1$.

(ii) Suppose such $V,\bar\pi,C',\rho$ exist.  First, $\Pi=\mathbf 1\otimes\bar\pi$
is bounded on the $V$-space, so $\bar\pi V<\infty$ and
$\kappa_1V\le(C'\rho+\bar\pi V)\,V$; in particular $\kappa_1V(x_0)<\infty$.
Second, for $y\in\mathcal Y$ and $y_\varepsilon$ obtained by carving a
sub-arc of mass $\varepsilon$ from its environment curve, the argument
of (i) at $y$ gives $\kappa_1^kF_i(y_\varepsilon)\ge1-2\varepsilon T-\Pr_y[\tau_k>T]$
and $\bar\pi F_i=\lim_k\kappa_1^kF_i(x_0)=0$, so with $T=T_k:=4Ck^p$
and $\varepsilon\le\varepsilon_k:=1/(8T_k)$ (Markov and $(\mathrm M_p)$),
\[
V(y_\varepsilon)\;\ge\;\frac{1-\tfrac14-\tfrac14}{C'\rho^k}\;=\;\frac1{2C'\rho^k}
\qquad(y\in\mathcal Y,\ \varepsilon\le\varepsilon_k).
\]
Third, from $x_0$ the first step carves a piece of mass in
$d\varepsilon$ from the environment curve with probability
$\ge\tfrac12\,d\varepsilon$ (pick that curve twice, $\tfrac14$;
$|t_1-t_2|/\tfrac12$ has density $\ge2$ near $0$) and the second step
is a flip landing in $\mathcal Y$ with probability $\ge\tfrac18\cdot\tfrac34$;
the landing state is then some $y_\varepsilon$, $y\in\mathcal Y$.  So
\[
\kappa_1V(x_0)\;\ge\;\frac{3}{64}\sum_{k\ge1}\frac{\varepsilon_k-\varepsilon_{k+1}}{2C'\rho^k}
\;\ge\;c\sum_k\frac{k^{-p-1}}{\rho^{k}}\;=\;\infty ,
\]
contradicting the first step.
\end{proof}

\emph{Status: PROVED ((ii) under $(\mathrm M_p)$, which (R1) along the
chain implies with $p=1$ and which is measured: $\E[\tau_k]\approx6k$).
For the chain killed on $\Omega_K$ the same argument runs up to
$k\le c\,3^{K/(3p)}$, so any $V$-uniform ergodicity of the killed chain
has $(1-\rho)^{-1}\log(C'V(x_0)\bar\pi V)\ge c\,3^{K/(3p)}$: exponential
in $K$, i.e.\ fatal in the budget of the session-11 handoff.  This
settles, negatively, the sentence of \S\ref{sec:markovrenewal} that
asked for a weighted space ``in which the post-flip chain is
$V$-geometrically ergodic in the base coordinates'': there is none.
It does \emph{not} touch (C1) or (C2) themselves, whose null vectors
would have to sit near $-1$, whereas the dust modes sit near $+1$.}

\subsubsection{The drift for \texorpdfstring{$\lambda$}{lambda} fails
on the full state space (TODO~12a)}\label{sec:s12drift}

\begin{proposition}[no Lyapunov drift for $(\sum r_j^4)^{-\theta}$]
\label{prop:nodrift}
Let $V=(\sum_jr_j^4)^{-\theta}$ with $0<\theta\le\tfrac12$, on the full
(unkilled) state space, and assume $\tau_k<\infty$ a.s.\ from every
state.  There are no $\varrho<1$ and $b<\infty$ with
$\kappa_1V\le\varrho V+b$.
\end{proposition}

\begin{proof}
Suppose the inequality holds.  Iterating,
$\kappa_1^kV\le\varrho^kV+b/(1-\varrho)$ for all $k$.  Fix $k$ with
$\varrho^k\le\tfrac38\,17^{-\theta}$ and $m\in(0,\tfrac14)$.

\emph{The isolated process.}  Let $\bar x$ be the apart-sector state
with $a=m$, $b=m/10$ and one environment curve of mass $1-1.1m$, and
let $\mathrm P^{\rm iso}$ be the law of the chain from $\bar x$ modified
so that a chord with one point in the original environment curve (or in
any curve descended from it by splits and merges among such curves)
and one point elsewhere does nothing.  Under $\mathrm P^{\rm iso}$ the
``non-dust'' material of total mass $1.1m$ evolves as a split--merge
chain in its own right, slowed by the factor $(1.1m)^2$, with the
marked pair inside it; its flip times are a.s.\ finite, so there is
$T=T(m,k)$ with $\mathrm P^{\rm iso}[\tau_k>T]\le\tfrac18$.

\emph{The states $x_{m,N}$.}  Let $x_{m,N}$ be $\bar x$ with the
environment curve cut into $N$ equal pieces,
$N\ge16(T+1)^3\max(m^{-1},m^{-4/3})$.  Drive the chain from $x_{m,N}$
and the isolated process from $\bar x$ by the same points (the dust
pieces of $x_{m,N}$ are sub-arcs of the environment curve of $\bar x$;
dust--dust chords split or merge dust material in both and never
touch the non-dust material).  The two processes make the same non-dust moves
except at chords joining dust to non-dust, of which there are at most
$T$ up to time $T$; each such chord appends to one non-dust curve of the
$x_{m,N}$-chain a dust curve, i.e.\ a union of at most $T+1$ original pieces or
fragments of them, of mass $\le(T+1)/N$.  As in Lemma~\ref{lem:inert} (with the
appended material in the role of $D$, now inside a curve of one chain
and absent from the other), the non-dust configurations of the two
processes then agree up to the appended material, and their flip times
agree, until a point falls in appended material, which has probability
$\le2T(T+1)/N$ per step.  Hence
$\Pr_{x_{m,N}}[\tau_k>T]\le\tfrac18+2T^2(T+1)/N\le\tfrac14$.

\emph{The deterministic bound.}  Along the $x_{m,N}$-chain up to time
$T$, the non-dust mass is at most $1.1m+T(T+1)/N\le2m$ and every dust
curve has mass $\le(T+1)/N\le m^{4/3}$, so
\[
\sum_jr_j(\bar\Gamma_t)^4\;\le\;(2m)^4+\Bigl(\max_{\text{dust}}r_j\Bigr)^3
\sum_{\text{dust}}r_j\;\le\;16m^4+m^4=17m^4\qquad(t\le T),
\]
whence $\kappa_1^kV(x_{m,N})\ge\Pr_{x_{m,N}}[\tau_k\le T]\,(17m^4)^{-\theta}
\ge\tfrac34\,17^{-\theta}m^{-4\theta}$, while $V(x_{m,N})\le m^{-4\theta}$
because the $u$-curve alone contributes $m^4$.  The drift bound gives
$\tfrac34\,17^{-\theta}m^{-4\theta}\le\tfrac38\,17^{-\theta}m^{-4\theta}+b/(1-\varrho)$,
i.e.\ $m^{-4\theta}\le\tfrac83\,17^{\theta}b/(1-\varrho)$ for every
$m\in(0,\tfrac14)$: false.
\end{proof}

\emph{Status: PROVED.  This is the fallback the brief provided for: the
obstruction is the family ``all non-dust mass $\asymp m$, the rest in
ultra-fine mass-carrying dust'', which is of vanishing probability
under the chain from one curve (Corollary~\ref{cor:horizon}) but is a
legitimate state, and on it $\lambda$ cannot recover within one
inter-flip period because there is no macroscopic partner to merge
with.  On $\Omega_K$, $V\le(2K)^{4\theta}$ is bounded and the drift is
vacuous.  The killing at $B=K$ in Theorem~\ref{thm:idleloop} therefore
stays, and it is harmless: it costs $(9d+1)^{-\gamma}$ over the horizon.
The second purpose of TODO~12a --- supplying the base half of the
minorisation --- is void by Proposition~\ref{prop:nomino}.  A quick
numerical check of the mechanism (\texttt{s12\_driftfail}): from
$(0.3,0.03;\ 3\times10^4\text{ equal pieces})$, $\E[(\sum r^4(\bar\Gamma_\tau)/\sum r^4(x))^{-\theta}]=1.19$
$(\theta=0.1)$, $1.71$ $(\theta=0.3)$, against $0.91$, $0.77$ with a
single environment curve of the same mass.}

\subsubsection{TODO~12c: the stated criterion does not give (C2)}
\label{sec:s12odd}

The session-11 handoff proposed, as route 12c, a state-uniform lower
bound $c(K)$ on the probability that two consecutive inter-flip times
have opposite parity, ``which gives (C2) by a standard two-block
argument''.  It does not: if the inter-flip times were the
deterministic sequence $1,2,1,2,\dots$ then consecutive parities are
always opposite ($c=1$) while $(-1)^{\tau_k}$ is periodic with period
$4$ and does not decay.  What the two-block argument needs is the
\emph{operator} statement $\|\kappa_{-1}^{\,2}\|_\infty<1$, and that is
false exactly: $\kappa_{-1}(-1)^E=-(-1)^E$ (Proposition~\ref{prop:r2false}),
so $\|\kappa_{-1}^{\,m}\|_\infty=1$ for every $m$.  Any odd perturbation
that is exact returns a state with a different block-count parity,
which is visible; the only odd perturbations that are invisible are
dust insertions, whose weight at scale $\eta$ is $\asymp2\eta$ per step
and whose lifetime is $\asymp1/(2\eta)$, and one checks in one line
that they give a contraction $1-O(\eta)$ per flip and nothing
geometric.  Route 12c is closed as stated; the honest form of (C2) is
in \S\ref{sec:s12num} and \S\ref{sec:s12left}.

\subsubsection{The decay of the parity along the post-flip chain is
polynomial}\label{sec:s12num}

By Lemma~\ref{lem:parity}, $(-1)^{\tau_k}=(-1)^{B(\bar\Gamma_{\tau_k})-B(\bar x)}$
exactly, so
\begin{equation}\label{eq:akB}
a_k(\bar x):=\E_{\bar x}[(-1)^{\tau_k}]=(-1)^{B(\bar x)}\sum_{b\ge1}(-1)^b\,
\Pr_{\bar x}\bigl[B(\bar\Gamma_{\tau_k})=b\bigr],
\end{equation}
the parity bias of the block count at the $k$-th flip, and
$\kappa_1^{\,k}(-1)^E=(-1)^{E}(-1)^ka_k$, so (C2) is the statement that
\eqref{eq:akB} decays geometrically, uniformly in $\bar x$.

\emph{Measurement (\texttt{s12\_akdecay}, \texttt{s12\_akB}; fixed
deterministic starts, $2\times10^7$ independent trajectories each,
standard error $2.2\times10^{-4}$, no censoring).}  From the one-curve
start (arcs $(\tfrac12,\tfrac12)$):
\[
a_k=-0.4003,\ +0.1530,\ -0.0674,\ +0.0335,\ -0.0186,\ +0.0111,\ -0.0073,\
+0.0048,\ -0.0034,\ +0.0022,\ -0.0019,\ +0.0015\quad(k=1,\dots,12),
\]
so that $k^3|a_k|=2.51,2.48,2.49,2.20,2.54,2.54$ for $k=7,\dots,12$:
\[
a_k(\text{one curve})\;=\;(-1)^k\,(2.5\pm0.1)\,k^{-3}\qquad(7\le k\le12),
\]
with the successive ratios $|a_{k+1}/a_k|$ rising through
$0.38,0.44,0.50,0.56,0.60,0.66,0.66,0.70$ --- no geometric rate.  A
$10^8$-sample run to $k=20$ (standard error $10^{-4}$) gives
$|a_k|=0.00495,0.00349,0.00271,0.00195,0.00142,0.00115,0.00107,0.00095,
0.00091,0.00079,0.00055,0.00045,0.00035$ for $k=8,\dots,20$, all with
sign $(-1)^k$; $k^3|a_k|$ stays in $2.4$--$2.7$ up to $k=13$ and
fluctuates in $2.8$--$3.9$ (two to three standard errors) for
$k=14$--$20$; the least-squares exponent over $8\le k\le20$ is
$2.7\pm0.3$.  From
the three-curve apart start $(0.3,0.3;\{0.4\})$: $a_1=-0.004$,
then $-0.0606,+0.0385,-0.0219,+0.0132,-0.0086,+0.0060,-0.0041,+0.0031,
-0.0020,+0.0018,-0.0014$ ($k=2,\dots,12$), with $k^3|a_k|\to2.0$--$2.4$:
the same tail, sign $(-1)^{k+B_0+1}$.  Inserting a dust curve of mass
$10^{-4}$ into either start (mass taken from an environment curve or
from an arc) reproduces every $a_k$ to within two standard errors, as
Lemma~\ref{lem:inert} predicts.  From a $9$-curve start
$(0.3,0.3;\ 7\times0.057)$, $10^8$ samples: $a_k$ passes through zero
at $k=3$ and then \emph{grows} in absolute value to $0.00145$ at
$k=7$--$8$ before decaying, with the sign $(-1)^{k+1}$ at every
$k=4,\dots,20$ (seventeen consecutive alternations); $k^3|a_k|$ rises
from $0.5$ at $k=7$ to $2$--$3$ at $k=15$--$20$.  So the tail is
universal --- $C k^{-3}$ with $C\approx2.5$--$3$ --- and is reached
after a start-dependent transient, of about eight flips when the
environment starts with nine curves.  From a $17$-curve start
($5\times10^6$ samples) $|a_k|\le7\times10^{-4}$ for $3\le k\le16$,
the transient not yet over.  From the
burnt-in ``stationary'' start of session~11 (gap window
$[\tfrac18,\tfrac38)$, $5\times10^6$ samples), $a_k(\mu)$ is at the noise
floor $4.5\times10^{-4}$ for $k\ge5$ --- but that average carries the
factor $(-1)^{B(\bar x)}$ of \eqref{eq:akB} inside the mixture, which
cancels the tail; the sign-blind $\E_\mu[a_k(\bar x)^2]$ is at
$(4\pm5)\times10^{-4}$ for $k=5..8$, consistent with the fixed-start
values.  The decomposition of \eqref{eq:akB} by $b$ from the one-curve
start (\texttt{s12\_akB}) shows the mechanism: $\Pr[B_{\tau_k}=1]$
decays like $c/k$ with $c\approx0.5$ (the chain merges back to one curve with
probability $\asymp1/t$, and there $\tau_k$ is even), the law of
$B_{\tau_k}$ spreads only like $\log k$, and the alternating sum of a
lattice law that spreads logarithmically decays polynomially.

\emph{Status: MEASURED.  Conclusions.  (i) (C2) in the geometric form
of \S\ref{sec:s11left} is false on the full state space, and on
$\Omega_K$ it holds only with $1-\rho(K)\le e^{-cK}$ (the rate the
killing provides), which is outside the budget; the session-11 value
$\rho\approx0.3$ described the transient $k\le4$.  (ii) The polynomial
decay is a $+1$ phenomenon, not a $-1$ phenomenon: the generating
function $\sum_ka_kw^k$ is singular at $w=-1$, i.e.\ at the eigenvalue
$-1$ of $\kappa_{-1}$, which is the eigenvalue $+1$ of $\kappa_1$ (the
trivial one), approached sub-geometrically by the slow variable $B$.
Nothing in the measurement bears on (C1).  (iii) The decay is
sign-alternating in $k$ from every start tested, with the sign
$(-1)^{k+B_0+1}$; this is the statement that $B_{\tau_k}\equiv k+1$ is
slightly favoured, i.e.\ that the block count at the $k$-th flip
remembers that it started odd --- a memory that $\log$-growth of $B$
erases only polynomially.}

\subsubsection{What is left, restated}\label{sec:s12left}

\begin{proposition}[polynomial parity decay suffices at $z=-1$]
\label{prop:holder}
Fix $K$ and $S'_K$.  Suppose
\begin{itemize}[leftmargin=2em]
\item[(C2$''$)] there are $C,p$ and $\delta_0>0$ with
$\sup_{\bar x\in S'_K}\bigl|\E_{\bar x}[w^{\tau_k}(-1)^{E_{\tau_k}}]\bigr|\le Ck^{-p}$
for all $k\ge1$ and all $|w|\le1$ with $|1-w|\le\delta_0$;
\item[(M$_1$)] $\sup_{\bar x\in S'_K}\E_{\bar x}[\tau_k]\le C_1k$ for all $k$.
\end{itemize}
Then $G(z)=(I+\kappa_z)^{-1}(I-\kappa_z)\mathbf 1$, evaluated on $S'_K$,
is H\"older continuous of exponent $(p-1)/(p+1)$ on
$\{|z|\le1,\ |1+z|\le\delta_0\}$, with constant
$\le C_2\,(C+C_1)$.  In particular $p>13/5$ gives an exponent $>4/9$.
\end{proposition}

\begin{proof}
By Proposition~\ref{prop:repair}, for $z=-w$,
$G(-w)=\bar M(I-\kappa_w)^{-1}(I+\kappa_w)\bar M\mathbf 1
=\mathbf 1+2\bar M\sum_{k\ge1}g_k(w)$ with
$g_k(w)(\bar x)=\E_{\bar x}[w^{\tau_k}(-1)^{E_{\tau_k}}]$
(Neumann series, $\|\kappa_w\|\le|w|<1$; the boundary values are limits).
For $|w|,|w'|\le1$, $|w^n-w'^n|\le n|w-w'|$, so
$|g_k(w)-g_k(w')|\le\min\bigl(2Ck^{-p},\ C_1k|w-w'|\bigr)$ by (C2$''$)
and (M$_1$).  Summing with the cut at $k_*=|w-w'|^{-1/(p+1)}$:
$\sum_{k\le k_*}C_1k|w-w'|+\sum_{k>k_*}2Ck^{-p}
\le C_2(C+C_1)\,|w-w'|^{(p-1)/(p+1)}$.
\end{proof}

\emph{Status: PROVED as a reduction; the hypotheses are OPEN.  Two
comments.  (a) The exponent $(p-1)/(p+1)$ is crude: it bounds
$|g_k(w)-g_k(w')|$ by the trivial Lipschitz bound and never uses that
$g_k(1)=(-1)^{E}(-1)^ka_k$ has one sign for all $k$.  If the parity
memory factorises from the time weighting,
\[
\text{(C2$'''$)}\qquad
\E_{\bar x}[w^{\tau_k}(-1)^{\tau_k}]=\E_{\bar x}[w^{\tau_k}]\,a_k(\bar x)\,(1+O(|1-w|k)),
\]
then $\sum_kg_k(w)\approx(-1)^E\sum_kCk^{-p}\E[w^{\tau_k}]$, and with
$\E[w^{\tau_k}]\approx w^{6k}$ this is $C\,\mathrm{Li}_p(w^6)$ up to
smoother terms, whose modulus of continuity at $w=1$ is
$|1-w|^{\min(1,p-1)}$ (up to a logarithm at integer $p$): the budget
then needs only $p>13/9$.  (C2$'''$) was measured (\texttt{s12\_twist}, one-curve start,
$3\times10^7$ samples, $w=0.99,0.95,0.90$ real): the exact
factorisation fails --- the twisted parity $\E[w^{\tau_k}(-1)^{\tau_k}]$
exceeds $\E[w^{\tau_k}]\,a_k$ by a factor $1.2$ at $|1-w|k=0.08$ and
$2$--$4$ at $|1-w|k\approx0.4$--$0.8$ (short paths carry more parity
memory) --- but the uniform bound (C2$''$) holds with the same
constant, $|\E[w^{\tau_k}(-1)^{\tau_k}]|\le|a_k|$ at every $k\le16$ and
every $w$ tested.  More usefully, the sum itself was measured: with
$\Sigma(w):=\sum_k(-1)^kE_x[w^{\tau_k}(-1)^{\tau_k}]$ (all terms of one
sign), $\Sigma(1)-\Sigma(w)=0.026,\ 0.085,\ 0.143$ at
$|1-w|=0.01,0.05,0.10$ (tail beyond $k=16$ estimated from $2.5k^{-3}$),
i.e.\ a modulus of continuity $\asymp|1-w|^{0.75}$ over that range ---
above the crude $(p-1)/(p+1)\approx0.5$, below the $\min(1,p-1)$ of the
factorised heuristic, and far above the $4/9$ the budget needs.  The
same on the unit circle near $w=1$ is the first numerical task of
session~13.  With the measured $p=2.7\pm0.3$ the crude route is marginal
($p>13/5$ required); the measured modulus is not.  (b) (C2$''$) at $w=1$ is exactly the polynomial
(C2$'$): $\sup_{S'_K}|a_k|\le Ck^{-p}$.  Away from $w=1$ on the circle
Theorem~\ref{thm:idleloop} gives a geometric rate; (C2$''$) asks that
the polynomial rate survive a small complex twist near $w=1$, which is
a damping of long paths and should not spoil a parity cancellation that
is produced by the block count.  This is measurable and is the first
numerical task for session~13.}

\medskip
\noindent The open conditions are now:

\medskip
\noindent\textbf{(C1)}\quad$\|(I+\kappa_1)^{-1}\|_{L^\infty(\Omega_K)}\le C(K)$ ---
equivalently $\|(I-\kappa_{-1})^{-1}\|\le C(K)$, the summability of the
signed renewal series $\sum_k\kappa_{-1}^{\,k}G$ for every bounded $G$;
(C2$'$) is the case $G=\mathbf 1$.  Unchanged, not contradicted, and no
longer reachable by any Doeblin/ergodicity argument
(Proposition~\ref{prop:nomino}, Corollary~\ref{cor:noerg}).\\
\noindent\textbf{(C2$''$)}\quad as in Proposition~\ref{prop:holder},
with $p>13/5$ and $C=C(K)$ polynomial; measured $p\approx3$ at $w=1$.

\medskip
\emph{What a proof must look like.  (i) It must be dust-blind: by
Lemma~\ref{lem:inert}, $a_k(y_\varepsilon)=a_k(y)+O(\varepsilon\,\E_y\tau_k)$
while $(-1)^{E}$ changes sign, so any argument that treats $(-1)^E$ as
a function to be contracted in a sup-type norm on configurations is
doomed, and any argument that works must use a topology in which
$y_\varepsilon\to y$.  (ii) The natural candidate --- the transport
distance between mass configurations, under the coupling of
Lemma~\ref{lem:inert} --- is not contracted: when the dust is finally
hit, the two coupled configurations differ by a macroscopic merge
(``$r+r_0$'' against ``$r+\varepsilon$, $r_0-\varepsilon$''), so a
discrepancy of size $\varepsilon$ is amplified to size $O(1)$ rather
than resolved.  A contracting metric would have to compare coarse
\emph{laws} rather than configurations.  (iii) The environment's own
randomness is needed: an adversary allowed to choose the masses gained
by the marked pair at the merge steps (whose timing it cannot control)
can bias the parity of each inter-flip time toward a target by a fixed
amount, since a gain changes the flip hazard $2p_1p_2$ by a bounded
amount at a time of its choosing among the merge steps; so no
``uniformly over environments'' statement can be true, and the proof
must use that the true gains are size-biased picks from a partition
whose count parity is invisible to them.  (iv) Given
\eqref{eq:akB}, the cleanest target is a statement about the block
count alone: the parity of $B$ at the $k$-th flip equidistributes
polynomially.  Since the total block count from one curve has the
exact law of Theorem~\ref{thm:blockcount} at fixed times, and the flip
times are a renewal-like sampling of the time axis, this may be
provable from the exact formulas by an averaging argument, at least
for the annealed environment --- which is the form the downstream
reduction (Problem~\ref{prob:survival}, whose start is
$\mu_j$) actually consumes.}

%----------------------------------------------------------------------
\subsection{Session 13: audit of the reduction chain}\label{sec:s13}

Session 13 opened with the audit recommended in the session-11 health
note: nothing in the chain Theorem~\ref{thm:marked} $\to$
Problem~\ref{prob:annealed} $\to$ Theorem~\ref{thm:factor} $\to$
Problem~\ref{prob:mixing} $\to$ Problem~\ref{prob:survival} $\to$
(R1)+(C1)+(C2$''$) had been re-read since it was laid, and two carried
conditions had just been found false in consecutive sessions.  The
audit produced three corrections to the \emph{form} of the bottom of
the chain, none of which changes what is proved above it:
\begin{enumerate}[leftmargin=2.2em]
\item[(1)] the quantity the programme consumes is an $L^2$ average of
the quenched parity $\Phi_m$ over the start, not a supremum over
states (Proposition~\ref{prop:S2red}); the state-uniform target (S)
of \S\ref{sec:s12left} was introduced by the strong-Markov step of
the session-10 reduction, not demanded by anything above it, and the
inert-dust obstruction of \S\ref{sec:s12dust} does not touch the
$L^2$ form (Remark~\ref{rem:dustblind});
\item[(2)] the exponent budget carried since session 10
($\gamma>4/9$ at $z=-1$, hence $p>13/5$ in
Proposition~\ref{prop:holder}) is too demanding by a factor of
about five: the correct budget is $\gamma>1/11$, hence $p>6/5$
(Corollary~\ref{cor:budget}).  The measured $p\approx2.7$ and
modulus $0.75$ are far inside it, not marginal;
\item[(3)] $\tau_k<\infty$ a.s.\ from every state with $ab>0$, taken
for granted since session 10, has a four-line proof
(Lemma~\ref{lem:flipsrecur}).
\end{enumerate}
The occupation-time step and (M$_1$) are stated exactly and left open
(\S\ref{sec:s13occ}, \S\ref{sec:s13m1}).

\subsubsection{The reduction consumes an \texorpdfstring{$L^2$}{L2} average, not a supremum}
\label{sec:s13ltwo}

Recall (Corollary~\ref{cor:quotient}) that
$\psi(j,d)=\E_{\mu_j}[\Phi_d(\Gamma_0)^2]$ with
$\Phi_m(\Gamma)=\E_\Gamma[(-1)^{N_m}]$ the parity of the number of
flips in $m$ steps of the quotient chain, and that
Theorem~\ref{thm:factor} needs $\psi(j,d)\le Kd^{-2c}$ for
$j\le8d$ only.  Fix $K\ge3$ and
$S'_K=\{a,b\ge1/(16K)\}$ as in \S\ref{sec:markovrenewal}.

\begin{proposition}[$L^2$ form of the reduction to $S'_K$]\label{prop:S2red}
Let $d\ge1$, $j\le8d$, $0<\varepsilon\le1/(2K)$, and let
$T^*=T^*(\varepsilon,K,d)<d/2$ be any deterministic time.  Let
$T'=\inf\{t\ge0:\Gamma_t\in S'_K\}$ and let
$\nu=\nu_{j,\varepsilon,K,T^*}$ be the joint law of
$(\Gamma_{T'},T')$ under $\Pr_{\mu_j}(\,\cdot\,;\,g_0\ge\varepsilon,\
T'\le T^*)$, a sub-probability measure on $S'_K\times\{0,\dots,T^*\}$;
write $\nu_t=\nu(\cdot\times\{t\})$ and, when $\nu_t\ne0$, let
$\bar\nu_t=\nu_t/\nu_t(S'_K)$ be the \emph{time-$t$ entrance law}, a
probability measure on $S'_K$.  Then
\[
\E_{\mu_j}\bigl[\Phi_d(\Gamma_0)^2;\ g_0\ge\varepsilon\bigr]
\;\le\;2\int\Phi_{d-t}(x)^2\,\nu(dx,dt)
\;+\;2\,\Pr_{\mu_j}\bigl[T'>T^*;\ g_0\ge\varepsilon\bigr]
\;\le\;2\sup_{t\le T^*}\E_{\bar\nu_t}\bigl[\Phi_{d-t}^2\bigr]
\;+\;2\,\Pr_{\mu_j}\bigl[T'>T^*;\ g_0\ge\varepsilon\bigr].
\]
\end{proposition}

\begin{proof}
$T'$ is a stopping time and $T'\wedge T^*\le T^*<d$.  By the strong
Markov property at $T'\wedge T^*$ and the additivity of the flip
count, $\E[(-1)^{N_d-N_{T'}}\mid\mathcal F_{T'}]=\Phi_{d-T'}(\Gamma_{T'})$
on $\{T'\le T^*\}$, so
\[
\Phi_d(\Gamma_0)
=\E_{\Gamma_0}\bigl[(-1)^{N_{T'}}\Phi_{d-T'}(\Gamma_{T'});\ T'\le T^*\bigr]
+\E_{\Gamma_0}\bigl[(-1)^{N_d};\ T'>T^*\bigr]
=:A(\Gamma_0)+B(\Gamma_0).
\]
Then $|B|\le\Pr_{\Gamma_0}[T'>T^*]\le1$, so $B^2\le|B|$, and by
Cauchy--Schwarz
$A^2\le\E_{\Gamma_0}[\Phi_{d-T'}(\Gamma_{T'})^2;\,T'\le T^*]\cdot
\Pr_{\Gamma_0}[T'\le T^*]\le\E_{\Gamma_0}[\Phi_{d-T'}(\Gamma_{T'})^2;\,T'\le T^*]$.
Use $(A+B)^2\le2A^2+2B^2$ and integrate over $\Gamma_0\sim\mu_j$ on
$\{g_0\ge\varepsilon\}$; the first integral is, by definition of
$\nu$, $\int\Phi_{d-t}(x)^2\nu(dx,dt)=\sum_{t\le T^*}\nu_t(S'_K)\,
\E_{\bar\nu_t}[\Phi_{d-t}^2]$, and $\sum_t\nu_t(S'_K)\le1$.  (The
supremum over $t$ must stay inside the sum in this way: a bound in
terms of the time-marginal $\sum_t\nu_t$ alone would cost a factor
$T^*$.)
\end{proof}

\emph{Status: PROVED.  Compare the session-10 reduction paragraph at
the end of \S\ref{sec:markovrenewal}, which at the same point bounds
$|\Phi_d(\Gamma_0)|\le\E|\Phi_{d-T'}(\Gamma_{T'})|+\Pr[T'>T^*]$ and
then takes the supremum of $|\Phi_m|$ over $S'_K\cap H$.  The
supremum is where the state-uniform target (S) entered the
programme; nothing above this point asks for it.  What the programme
needs on $S'_K$ is therefore:}

\medskip
\noindent\textbf{(S$_2$)}\quad
$\displaystyle\E_{\bar\nu_t}\bigl[\Phi_{m}^2\bigr]\le C(K)\,m^{-2\gamma}$
\emph{for every $t\le T^*$ and $m=d-t\in[d/2,d]$, for the time-$t$
entrance laws $\bar\nu_t=\bar\nu_{j,\varepsilon,K,T^*;t}$ of $S'_K$,
$j\le8d$, with $K=K_{\gamma'}(9d)=O(\log d)$ and $C(K)$ polynomial
in $K$,}

\medskip
\noindent together with the entrance estimate
$\Pr_{\mu_j}[T'>T^*;\,g_0\ge\varepsilon]\le O(d^{-\gamma'})+O(\varepsilon K\log d)$
for $T^*=O(\varepsilon^{-1}K\log d)$, which is the content of the
session-10 reduction paragraph (Lemma~\ref{lem:srecharge} plus the
occupation-time step of \S\ref{sec:s13occ}).  The block-bound event
$H_{9d}$ of Corollary~\ref{cor:horizon} is also handled in $L^2$
without any state-by-state argument: with $H=\{B_r<K,\ r\le m\}$
for the chain from $x$, writing $\Phi_m=\Phi_m^H+\Phi_m^{H^c}$ with
$|\Phi_m^{H^c}(x)|\le\Pr_x[H^c]$, one has
$\E_{\bar\nu_t}[\Phi_m^2]\le2\E_{\bar\nu_t}[(\Phi_m^H)^2]+2\E_{\bar\nu_t}[\Pr_x[H^c]]$,
and $\sum_t\nu_t(S'_K)\,\E_{\bar\nu_t}[\Pr_x[H^c]]\le\Pr_{\mu_j}[H^c_{9d}]\le(9d+1)^{-\gamma'}$
because $\nu$ is a sub-law of the path measure from $\mu_j$, whose
clock reaches at most $j+d\le9d$.

\begin{remark}[what the dust obstruction does and does not touch]
\label{rem:dustblind}
(a) Proposition~\ref{prop:nomino} and Corollary~\ref{cor:noerg} are
statements about the supremum norm on $\Omega_K$: they show that no
Doeblin-type argument can give a bound on $\sup_{x\in S'_K}|\Phi_m(x)|$
(or on $\|\kappa_1^kF-\bar\pi F\|_\infty$ for any weight).  They say
nothing against (S$_2$), which is an average over laws $\bar\nu_t$
that are themselves marginals of the chain run from one block.
(b) Under the coupling of Lemma~\ref{lem:inert}, $\Phi_m$ is
Lipschitz in the dust: $|\Phi_m(y_\varepsilon)-\Phi_m(y)|\le4\varepsilon m$,
since the flip counts agree on the coupling event.  So $\Phi_m$ is a
continuous function in the coarse topology at scale $1/m$ --- the
topology in which $y_\varepsilon\to y$ --- and the dust of
Proposition~\ref{prop:nomino}, which lives at scales
$\varepsilon_i\ll1/T_i$, is invisible to the functions that (S$_2$)
is about.  This is the precise sense in which \S\ref{sec:s12left}(i)
asked for a ``dust-blind'' statement, and (S$_2$) is one.
(c) (S$_2$) is \emph{not} an annealed statement in the naive sense:
it bounds the second moment of the quenched parity, so a proof must
still show that $\Phi_m(x)$ is small for $\bar\nu_t$-typical $x$, with
the environment of $x$ frozen; only the \emph{uniformity} over $x$
has been removed.  The quantity $\E_{\bar\nu_t}[\Phi_m]$ (parity
averaged over the environment of the start) is not what is needed and
is not sufficient.
(d) The natural function space for (S$_2$) is $L^2(\bar\pi)$, in
which the chain is self-adjoint (\S\ref{sec:spectral}); the spectral
route of session 8, shelved in favour of the Markov-renewal route,
is the dust-blind route.  What it additionally needs is a comparison
of the entrance laws $\bar\nu_t$ with $\bar\pi$ on the functions
$\Phi_m^2$; this is the first sub-problem of \S\ref{sec:s13spectral}.
\end{remark}

\subsubsection{The exponent budget, recomputed}\label{sec:s13budget}

\begin{corollary}[the budget is $\gamma>1/11$, not $\gamma>4/9$]
\label{cor:budget}
Suppose (S$_2$) holds with exponent $\gamma\in(0,\tfrac12)$ and the
entrance estimate holds as stated after Proposition~\ref{prop:S2red}.
Then $\psi(j,d)\le C'(\log d)^{O(1)}\,d^{-2\gamma}$ for all $j\le8d$,
i.e.\ Problem~\ref{prob:mixing} holds with $c=\gamma$, and
Corollary~\ref{cor:psienough} closes the iid model as soon as
$\gamma>1/11$.  The same holds if (S$_2$) is replaced by the
state-uniform (S) with exponent $\gamma$.  In particular, in
Proposition~\ref{prop:holder} the crude route needs only $p>6/5$,
and (R1) needs only $C^{1,\gamma}$ regularity with $\gamma>1/11$,
i.e.\ moments $\E_x[\tau^{1+\gamma}]$ with $\gamma>1/11$.
\end{corollary}

\begin{proof}
Set $\varepsilon:=d^{-1/2-\gamma}$ and $T^*:=C\varepsilon^{-1}K\log d$
with $K=K_{\gamma'}(9d)$, $\gamma'\ge2\gamma$; then
$T^*=O(d^{1/2+\gamma}\log^2d)<d/2$ for $d$ large because
$\gamma<\tfrac12$, and $\varepsilon\le1/(2K)$.  By
Proposition~\ref{prop:shorthorizon} ($j\le8d$),
$\Pr_{\mu_j}[g_0<\varepsilon]\le4\varepsilon^2j\le32\varepsilon^2d=32d^{-2\gamma}$.
By Proposition~\ref{prop:S2red} and (S$_2$) (with $m=d-t\ge d/2$),
$\E_{\mu_j}[\Phi_d^2;g_0\ge\varepsilon]\le2C(K)(d/2)^{-2\gamma}
+2\bigl(O(d^{-\gamma'})+O(\varepsilon K\log d)\bigr)$, and
$\varepsilon K\log d=O(d^{-1/2-\gamma}\log^2d)=o(d^{-2\gamma})$ since
$\gamma<\tfrac12$.  Add, and use $|\Phi_d|\le1$ on $\{g_0<\varepsilon\}$.
The budget: Corollary~\ref{cor:psienough} needs $c>1/11$.
For Proposition~\ref{prop:holder}: $(p-1)/(p+1)>1/11\iff p>6/5$.
\end{proof}

\emph{Status: PROVED (arithmetic).  Where the factor was lost.  The
session-10 reduction paragraph ends by bounding $|\Phi_d(\Gamma_0)|$,
not $\Phi_d(\Gamma_0)^2$, by $C(\varepsilon d)^{-2\alpha}$ --- the
form of Problem~\ref{prob:survival}, which is a bound on the
\emph{second moment} --- and then replaces $(\varepsilon d)^{-2\alpha}$
by $d^{-2\alpha}$.  The first step costs a factor $2$ in the exponent,
the second a further factor $2(1+\alpha)=22/9$; together they turn
the true requirement $\gamma>1/11$ into $\gamma>2\alpha>4/9$.
Problem~\ref{prob:survival} and Corollary~\ref{cor:kappafree} are
correct as stated and, with the correct accounting, lossless: from
$\Phi_d^2\le C(K)d^{-2\gamma}$ after the entrance one gets
Problem~\ref{prob:survival} with
$(\varepsilon d)^{-2\alpha}=d^{-\alpha/(1+\alpha)}$ at the optimising
$\varepsilon$ for $\alpha=2\gamma/(1-2\gamma)$, and then
$c'=\alpha/(2(1+\alpha))=\gamma$ exactly.  The loss was entirely in
the two steps just named.  Every downstream statement that quotes
``$\gamma>4/9$'' (\S\ref{sec:markovrenewal} (R1) discussion,
\S\ref{sec:s11left}, Proposition~\ref{prop:holder} and
\S\ref{sec:s12left}, the handoff documents of sessions 10--12) should
be read with $1/11$ in its place.  Consequences: the ``marginal''
verdict on the crude route of Proposition~\ref{prop:holder}
($p>13/5$ against a measured $2.7\pm0.3$) is withdrawn --- the crude
route needs $p>1.2$; (R1) needs barely more than a first moment of
$\tau$ along the chain (the measured tail is $t^{-2}$); and
Proposition~\ref{prop:lower}'s sharp rate $c=1/2$ exceeds the budget
by a factor $5.5$, as \S\ref{sec:factor} already noted for
$\bar\psi$.  The programme has a wide margin everywhere except in
the one place where it has no proof.}

The budget chain, for reference (each line is proved above it):
\[
\begin{array}{lll}
\toprule
\text{statement} & \text{exponent} & \text{needs}\\
\midrule
\E|\mathrm{bias}|=o(1/\log n),\ k\ge6\log^{11}n & k^{-c} & c>1/11\\
\psi(j,d)\le Kd^{-2c},\ j\le8d\ \text{(Thm.~\ref{thm:factor})} & c & c>1/11\\
\text{(S)}/\text{(S}_2)\ \text{on }S'_K\ \text{(Cor.~\ref{cor:budget})} & \gamma=c & \gamma>1/11\\
G\in C^{0,\gamma}\ \text{at }z=-1\ \text{(Prop.~\ref{prop:condecay})} & \gamma & \gamma>1/11\\
\text{(C2$''$) via Prop.~\ref{prop:holder}, crude} & (p-1)/(p+1) & p>6/5\\
\text{(C2$''$) via Prop.~\ref{prop:holder}, factorised} & \min(1,p-1) & p>12/11\\
\text{measured (\S\ref{sec:s12num})} & p=2.7\pm0.3,\ \text{modulus }0.75 & \\
\bottomrule
\end{array}
\]

\subsubsection{Flips recur: \texorpdfstring{$\tau_k<\infty$}{tau\_k<infty} a.s.}\label{sec:s13tau}

\begin{lemma}[flips recur]\label{lem:flipsrecur}
From every state of the quotient chain with $a_0b_0>0$,
$\sum_t2a_tb_t=\infty$ a.s.; consequently flips occur infinitely
often a.s., and $\tau_k<\infty$ a.s.\ for every $k$.
\end{lemma}

\begin{proof}
Write $m_t=\min(a_t,b_t)$, $M_t=\max(a_t,b_t)$ and
$p_t=2a_tb_t=2m_tM_t$, the conditional probability of a flip at step
$t+1$.  By Lemma~\ref{lem:splitmerge} (read on the quotient,
Corollary~\ref{cor:quotient}), at a non-flip step a marked coordinate
either stays fixed, increases (a merge), or is multiplied by
$\sqrt U$, $U\sim U(0,1)$ (a shrink: both cuts in that coordinate's
material, probability equal to the square of the coordinate).  Hence
$m_t$ is non-decreasing between the steps at which it strictly
decreases, and at a non-flip step $m$ can decrease in exactly two
ways: a shrink of the smaller coordinate, probability $m_t^2$; or a
shrink of the larger coordinate below $m_t$, probability
$M_t^2\cdot\Pr[M_t\sqrt U<m_t]=M_t^2(m_t/M_t)^2=m_t^2$.  So the event
$E_{t+1}=\{\text{step }t{+}1\text{ is not a flip and }m_{t+1}<m_t\}$
has $\Pr[E_{t+1}\mid\mathcal F_t]\le2m_t^2\le2m_tM_t=p_t$.

Let $\Omega_0=\{\sum_tp_t<\infty\}$.  By L\'evy's conditional
Borel--Cantelli lemma, on $\Omega_0$ almost surely only finitely many
flips occur and only finitely many $E_t$ occur.  All marked
coordinates stay strictly positive a.s.\ ($U>0$, flips land in the
open interval), so after the last flip and the last $E_t$ the process
$m_t$ is non-decreasing and bounded below by some $m_*>0$; then
$p_t\ge2m_*^2$ for all large $t$ and $\sum_tp_t=\infty$,
contradicting $\Omega_0$.  Hence $\Pr[\Omega_0]=0$, i.e.\
$\sum_tp_t=\infty$ a.s., and L\'evy's lemma in the other direction
gives flips infinitely often a.s.
\end{proof}

\emph{Status: PROVED.  This is the hypothesis of
Propositions~\ref{prop:r2false}, \ref{prop:nomino}, \ref{prop:nodrift}
and of the remark after Proposition~\ref{prop:genfun}; it needs
nothing about the environment --- in particular it holds with an
all-dust environment --- because the only mechanism that decreases the
smaller marked coordinate has probability at most the flip
probability.  It gives no rate; rates are (R1).}

\subsubsection{The occupation-time step, stated}\label{sec:s13occ}

The session-10 reduction (end of \S\ref{sec:markovrenewal}) and the
entrance estimate after Proposition~\ref{prop:S2red} use, without
proof, the following statement.

\begin{lemma}[occupation time of the marked total; OPEN]\label{lem:occupation}
Let $K\ge3$ and let $H=H_T$ be the block-bound event with bound $K$.
There are $c_K,C_K>0$ such that for every start with
$s_0\ge1/(2K)$, every $\varepsilon\le1/(8K)$ and every $T\ge1$,
\[
\Pr\Bigl[\#\{t\le T:\ s_t<\tfrac1{4K}\}>\tfrac T2\Bigr]
\;\le\;C_K\,e^{-c_K\varepsilon T/K}\;+\;2\varepsilon^2T\;+\;\Pr[H^c].
\]
\end{lemma}

\emph{Status: OPEN; believed standard.  What it is used for: in the
entrance estimate, after $s$ first reaches $1/(2K)$ (within
$T_1=O(\varepsilon^{-1}K\log d)$ steps by Lemma~\ref{lem:srecharge}),
one needs, on at least half of the next $T_2=O(\varepsilon^{-1}K\log d)$
steps, $s_t\ge1/(4K)$ and $\min(a_t,b_t)\ge\varepsilon$ (the second
from the flux bound, cost $4\varepsilon^2T_2$); on those steps a flip
has probability $\ge2\varepsilon(s_t-\varepsilon)\ge\varepsilon/(4K)$
and lands in $S'_K$ with conditional probability $\ge\tfrac12$
(the trapezoid of Lemma~\ref{lem:recharge}(ii) is symmetric and
unimodal), so $T'\le T_1+T_2$ except with probability
$O(d^{-\gamma'})+O(\varepsilon K\log d)$.  What fails if it is false:
the entrance estimate, hence the reduction of (S$_2$) to the
log-macroscopic set $S'_K$; (S$_2$) would then have to be proved for
the laws $\mu_j$ restricted to $\{g_0\ge\varepsilon\}$ with
$\varepsilon=d^{-1/2-\gamma}$, i.e.\ with the trap accounting inside
the $L^2$ statement.  Nothing above Problem~\ref{prob:survival}
depends on it.  Intended proof: Lemma~\ref{lem:logdrift} gives
$\E[\Delta\log s\mid\Gamma_t]\ge c\,s_t/K$ on $H\cap\{s_t\le1/(2K)\}$,
with down-jumps those of $\tfrac12\log U$; so for small $\theta>0$,
$V=s^{-\theta}$ is a supermartingale on $\{s\le1/(2K)\}\cap H$ with
drift $\le-\theta c's^{1-\theta}/K$, and the occupation time of
$\{\varepsilon\le s<1/(4K)\}$ over $T$ steps is at most
$\varepsilon^{-(1-\theta)}$ times the accumulated drift; the time
below $\varepsilon$ is priced by Lemma~\ref{lem:strap}; excursions
above $1/(2K)$ are handled by stopping $V$ at $1/(2K)$.  The missing
piece is only the bookkeeping of the re-entries into
$\{s<1/(2K)\}$ from above, each of which starts $V$ at
$\le(2K)^\theta\cdot2^\theta$ (a single shrink from $\ge1/(2K)$
lands at $\ge s\sqrt U$, and $\Pr[\sqrt U\le\tfrac12]=\tfrac14$) and
of which there are at most $O(T/K^2)$ in expectation (shrinks have
probability $\le s^2$).  We have not written it out.}

\subsubsection{What (M\texorpdfstring{$_1$}{1}) needs}\label{sec:s13m1}

Proposition~\ref{prop:holder} assumes (M$_1$):
$\sup_{\bar x\in S'_K}\E_{\bar x}[\tau_k]\le C_1k$.  By the strong
Markov property at the flip times,
$\E_{\bar x}[\tau_k]=\sum_{i<k}\E_{\bar x}\bigl[\E_{\bar\Gamma_{\tau_i}}[\tau]\bigr]$,
so (M$_1$) follows from a bound on the mean first-flip time that is
uniform \emph{on average over post-flip states}.  The post-flip state
has marked total $s$ (preserved by the flip) and landing gap $g'$
with $\Pr[g'\le\eta]\le\eta^2/(ab)$ (Lemma~\ref{lem:recharge}(ii)),
so $\E[1/g']\le\int_0^{s}\eta^{-2}\Pr[g'\le\eta]\,d\eta\cdot
\text{(const)}<\infty$ is automatic from a macroscopic pre-flip
state; what is needed is therefore
\[
\text{(R1$'$)}\qquad
\E_{\bar x}[\tau]\;\le\;C(K)\,\bigl(1+1/g(\bar x)\bigr)\quad
\text{on }H,\ \text{uniformly over the environment},
\]
together with control of the pre-flip state at each flip (so that
$1/(ab)$ in the landing bound is $O(K^2)$ on average), which is the
recurrence of the marked pair to $S'_K$ --- the same recharge
machinery as the entrance estimate.  (R1$'$) is the first-moment
part of TODO~11a (which asks for the tail $\Pr_x[\tau>t]\le C_Kt^{-2}$,
measured in \S\ref{sec:parity}(c)); by Corollary~\ref{cor:budget},
(R1) itself now needs only a $(1+\gamma)$-th moment with
$\gamma>1/11$, so a tail bound $t^{-1-\delta}$ for any $\delta>1/11$
would suffice for both.  (M$_p$) in Corollary~\ref{cor:noerg}(ii) is
used only in a negative result and needs no further attention.

\subsubsection{Where this leaves the programme}\label{sec:s13spectral}

Items (I)--(III) of the overview are unchanged in substance, but (I)
is now: prove (S$_2$) with any $\gamma>1/11$, for the time-$t$
entrance laws $\bar\nu_t$ of $S'_K$, with $C(K)$ polynomial.  Two routes remain, neither of
which is Doeblin-type:
\begin{itemize}[leftmargin=2em]
\item[(i)] \emph{The spectral route} (\S\ref{sec:spectral}): the
quotient chain is reversible for $\bar\pi$; for the odd function
$f=\varepsilon F$ the Dirichlet form has the massive flip term, and
polynomial decay of $\psi_{\bar\pi}(d)=\|P^d\varepsilon\|^2_{L^2(\pi)}$
with exponent $\beta$ is equivalent to the spectral measure of
$\varepsilon$ having mass $O(\delta^\beta)$ within $\delta$ of
$\pm1$.  A sufficient functional inequality is
\[
\text{(H$_\beta$)}\qquad
\bigl|\E_{\bar\pi}[F]\bigr|^2\;\le\;C\,\mathcal E_{\mathrm{odd}}(F)^{\beta}\,\|F\|_{L^2(\bar\pi)}^{2(1-\beta)}
\qquad\text{for all }F\in L^2(\bar\pi),
\]
with $\mathcal E_{\mathrm{odd}}(F)=\tfrac12\E_{\bar\pi}[(F-F')^2;\text{non-flip}]
+\tfrac12\E_{\bar\pi}[(F+F')^2;\text{flip}]$: applied to
$F=\Pi_\delta\varepsilon$ (spectral projection onto
$[1-\delta,1]$) it gives $\nu_\varepsilon([1-\delta,1])\le C\delta^\beta$.
By Proposition~\ref{prop:lower}, $\beta\le1$; the budget needs only
$\beta>2/11$.  (The limiting case $\beta=1$ is a Hardy inequality
in the gap coordinate and is false by a logarithm: test functions
$F=g^{-1}\log^{-2}(1/g)\mathbf 1\{g\le\delta\}$ have
$|\E F|^2/\mathcal E_{\mathrm{odd}}(F)\asymp\log(1/\delta)$ --- consistent
with $\psi_{\bar\pi}\asymp1/d$, since linear spectral density gives
$\int d\nu_\varepsilon/(1-\lambda)=\infty$.)  The periodicity
caveat of \S\ref{sec:spectral} applies: the same inequality is
needed for $(-1)^BF$, i.e.\ near $\lambda=-1$.  What the spectral
route does \emph{not} give by itself is (S$_2$) for the entrance
laws $\bar\nu_t$: these are singular with respect to $\bar\pi$
(finitely many blocks against infinitely many), and the comparison
must be made at a coarse scale --- the blocks of mass $\gtrsim1/(m\log m)$,
which are the only ones $\Phi_m$ sees (Remark~\ref{rem:dustblind}(b))
--- where $\bar\nu_t$ has had time of order $m\log m\gg$ the
relaxation time $1/(2\eta)$ of each scale $\eta$ in that range.
This start-comparison is the first concrete sub-problem.
\item[(ii)] \emph{The exact-law route} of \S\ref{sec:s12left}(iv):
polynomial decay of the block-count parity at flip times from the
exact fixed-time laws, from the one-curve start, extended to the
$\mu_j$ starts.  Its target is now $p>6/5$ in
Proposition~\ref{prop:holder}, i.e.\ $a_k=O(k^{-6/5})$, far weaker
than the measured $k^{-3}$.  Since $B(\bar\Gamma_{\tau_k})$ spreads
only like $\log k$, the decay cannot come from the width of its law;
it comes from the smoothness of the law of $\tau_k$ on the time axis
(the parity of $B$ at a \emph{fixed} time is deterministic,
Lemma~\ref{lem:parity}), and quantifying that smoothness is exactly
what (ii) must do.
\end{itemize}
Both routes consume the same two inputs that the audit has left open
--- the occupation-time lemma and (R1$'$) --- and both are now aimed at
a target with a fivefold margin.

